% Môn: Vật lý
% Lớp: 11
% Chương: Chương II. Sóng
% Bài: L11C2 Bài 15. Thực hành: Đo tốc độ truyền âm
% Số câu: 162
% App đọc file này trực tiếp — sửa rồi Commit trên GitHub, bấm Đồng bộ GitHub .tex

% L11C2 Bài 15. Thực hành: Đo tốc độ truyền âm

% ===== Câu 1 =====
% ID: L11C2B15-01-DS
% Mức: TH
\dangbt{Lắp đặt và sử dụng ống cộng hưởng}
\begin{ex}
% ID: L11C2B15-01-DS
% Mức: TH
Xét các phát biểu sau trong dạng “Lắp đặt và sử dụng ống cộng hưởng”.
\choiceTF
{\True Mục đích chính của thí nghiệm dùng ống cộng hưởng là xác định — phát biểu đúng là: tốc độ truyền âm trong không khí.}
{Mục đích chính của thí nghiệm dùng ống cộng hưởng là xác định — phát biểu cường độ dòng điện là đúng.}
{\True Dụng cụ phát âm có tần số xác định trong thí nghiệm thường là — phát biểu đúng là: âm thoa hoặc loa nối máy phát âm tần.}
{Dụng cụ phát âm có tần số xác định trong thí nghiệm thường là — phát biểu lực kế là đúng.}
\loigiai{
\begin{itemchoice}
\item (Đúng) Mục đích chính của thí nghiệm dùng ống cộng hưởng là xác định — phát biểu đúng là: tốc độ truyền âm trong không khí. (Vì): Từ tần số đã biết và bước sóng đo qua cộng hưởng, ta tính tốc độ truyền âm. b) Sai: phương án nêu không phù hợp. c) Đúng: Âm thoa hoặc loa tạo âm có tần số biết trước, cần cho phép đo. d) Sai: phương án nêu không phù hợp
\item (Sai) Mục đích chính của thí nghiệm dùng ống cộng hưởng là xác định — phát biểu cường độ dòng điện là đúng.
\item (Đúng) Dụng cụ phát âm có tần số xác định trong thí nghiệm thường là — phát biểu đúng là: âm thoa hoặc loa nối máy phát âm tần.
\item (Sai) Dụng cụ phát âm có tần số xác định trong thí nghiệm thường là — phát biểu lực kế là đúng.
\end{itemchoice}
}
\end{ex}

% ===== Câu 2 =====
% ID: L11C2B15-01-TL
% Mức: TH
\begin{ex}
% ID: L11C2B15-01-TL
% Mức: TH
Trong dạng “Lắp đặt và sử dụng ống cộng hưởng”, Nêu các bước chính đo tốc độ truyền âm bằng ống cộng hưởng.
\choiceTF

\loigiai{
Phát âm có tần số f đã biết; thay đổi chiều dài cột khí; ghi hai vị trí cộng hưởng liên tiếp; tính $\lambda=2|L_2-L_1|$ rồi tính $v=f\lambda$.
}
\end{ex}

% ===== Câu 3 =====
% ID: L11C2B15-01-TLN
% Mức: TH
\begin{ex}
% ID: L11C2B15-01-TLN
% Mức: TH
Trong dạng “Lắp đặt và sử dụng ống cộng hưởng”, Hai vị trí cộng hưởng liên tiếp cách nhau $34\,\text{cm}$. Tính bước sóng theo cm.
\shortans{68}
\loigiai{
$\lambda=2d=2\cdot34=68$ cm.
}
\end{ex}

% ===== Câu 4 =====
% ID: L11C2B15-01-TN
% Mức: NB
\begin{ex}
% ID: L11C2B15-01-TN
% Mức: NB
Trong dạng “Lắp đặt và sử dụng ống cộng hưởng”, Mục đích chính của thí nghiệm dùng ống cộng hưởng là xác định
\choice
{\True tốc độ truyền âm trong không khí}
{cường độ dòng điện}
{khối lượng riêng của nước}
{nhiệt dung riêng của không khí}
\loigiai{
Đáp án A. Từ tần số đã biết và bước sóng đo qua cộng hưởng, ta tính tốc độ truyền âm.
}
\end{ex}

% ===== Câu 5 =====
% ID: L11C2B15-02-DS
% Mức: TH
\begin{ex}
% ID: L11C2B15-02-DS
% Mức: TH
Xét các phát biểu sau trong dạng “Lắp đặt và sử dụng ống cộng hưởng”.
\choiceTF
{\True Khoảng cách giữa hai vị trí cộng hưởng liên tiếp trong ống một đầu kín xấp xỉ — phát biểu đúng là: $\lambda/2$.}
{Khoảng cách giữa hai vị trí cộng hưởng liên tiếp trong ống một đầu kín xấp xỉ — phát biểu $\lambda$ là đúng.}
{\True Biết khoảng cách hai vị trí cộng hưởng liên tiếp là d, bước sóng bằng — phát biểu đúng là: $2d$.}
{Biết khoảng cách hai vị trí cộng hưởng liên tiếp là d, bước sóng bằng — phát biểu $d$ là đúng.}
\loigiai{
\begin{itemchoice}
\item (Đúng) Khoảng cách giữa hai vị trí cộng hưởng liên tiếp trong ống một đầu kín xấp xỉ — phát biểu đúng là: $\lambda/2$. (Vì): Hai trạng thái cộng hưởng liên tiếp chênh nhau nửa bước sóng. b) Sai: phương án nêu không phù hợp. c) Đúng: Vì $d=\lambda/2$ nên $\lambda=2d$. d) Sai: phương án nêu không phù hợp
\item (Sai) Khoảng cách giữa hai vị trí cộng hưởng liên tiếp trong ống một đầu kín xấp xỉ — phát biểu $\lambda$ là đúng.
\item (Đúng) Biết khoảng cách hai vị trí cộng hưởng liên tiếp là d, bước sóng bằng — phát biểu đúng là: $2d$.
\item (Sai) Biết khoảng cách hai vị trí cộng hưởng liên tiếp là d, bước sóng bằng — phát biểu $d$ là đúng.
\end{itemchoice}
}
\end{ex}

% ===== Câu 6 =====
% ID: L11C2B15-02-TL
% Mức: TH
\begin{ex}
% ID: L11C2B15-02-TL
% Mức: TH
Trong dạng “Lắp đặt và sử dụng ống cộng hưởng”, Giải thích vì sao dùng hai vị trí cộng hưởng liên tiếp giúp xác định bước sóng.
\choiceTF

\loigiai{
Trong ống một đầu kín, các chiều dài cộng hưởng kế tiếp chênh nhau $\lambda/2$, nên đo hiệu hai vị trí loại được phần hiệu chỉnh đầu ống và cho $\lambda=2\Delta L$.
}
\end{ex}

% ===== Câu 7 =====
% ID: L11C2B15-02-TLN
% Mức: TH
\begin{ex}
% ID: L11C2B15-02-TLN
% Mức: TH
Trong dạng “Lắp đặt và sử dụng ống cộng hưởng”, Âm có tần số $500\,\text{Hz}$, bước sóng $0,68\,\text{m}$. Tính tốc độ theo m/s.
\shortans{340}
\loigiai{
$v=f\lambda=500\cdot0,68=340$ m/s.
}
\end{ex}

% ===== Câu 8 =====
% ID: L11C2B15-02-TN
% Mức: NB
\begin{ex}
% ID: L11C2B15-02-TN
% Mức: NB
Trong dạng “Lắp đặt và sử dụng ống cộng hưởng”, Dụng cụ phát âm có tần số xác định trong thí nghiệm thường là
\choice
{\True âm thoa hoặc loa nối máy phát âm tần}
{nhiệt kế}
{lực kế}
{cân điện tử}
\loigiai{
Đáp án A. Âm thoa hoặc loa tạo âm có tần số biết trước, cần cho phép đo.
}
\end{ex}

% ===== Câu 9 =====
% ID: L11C2B15-03-DS
% Mức: VD
\begin{ex}
% ID: L11C2B15-03-DS
% Mức: VD
Xét các phát biểu sau trong dạng “Lắp đặt và sử dụng ống cộng hưởng”.
\choiceTF
{\True Công thức tính tốc độ truyền âm là — phát biểu đúng là: $v=f\lambda$.}
{Công thức tính tốc độ truyền âm là — phát biểu $v=f/\lambda$ là đúng.}
{\True Để giảm sai số khi xác định khoảng cách cộng hưởng nên — phát biểu đúng là: đo nhiều cặp vị trí và lấy trung bình.}
{Để giảm sai số khi xác định khoảng cách cộng hưởng nên — phát biểu thay đổi tần số liên tục trong mỗi lần đo là đúng.}
\loigiai{
\begin{itemchoice}
\item (Đúng) Công thức tính tốc độ truyền âm là — phát biểu đúng là: $v=f\lambda$. (Vì): Tốc độ truyền sóng thỏa $v=f\lambda$. b) Sai: phương án nêu không phù hợp. c) Đúng: Đo lặp lại và lấy trung bình làm giảm ảnh hưởng sai số ngẫu nhiên. d) Sai: phương án nêu không phù hợp
\item (Sai) Công thức tính tốc độ truyền âm là — phát biểu $v=f/\lambda$ là đúng.
\item (Đúng) Để giảm sai số khi xác định khoảng cách cộng hưởng nên — phát biểu đúng là: đo nhiều cặp vị trí và lấy trung bình.
\item (Sai) Để giảm sai số khi xác định khoảng cách cộng hưởng nên — phát biểu thay đổi tần số liên tục trong mỗi lần đo là đúng.
\end{itemchoice}
}
\end{ex}

% ===== Câu 10 =====
% ID: L11C2B15-03-TL
% Mức: VD
\begin{ex}
% ID: L11C2B15-03-TL
% Mức: VD
Trong dạng “Lắp đặt và sử dụng ống cộng hưởng”, Với $f=400$ Hz, hai vị trí cộng hưởng là $18\,\text{cm}$ và $60\,\text{cm}$. Tính tốc độ âm.
\choiceTF

\loigiai{
$d=60-18=42\,\mathrm{cm}=0,42$ m; $\lambda=0,84$ m; $v=400\cdot0,84=336$ m/s.
}
\end{ex}

% ===== Câu 11 =====
% ID: L11C2B15-03-TLN
% Mức: VD
\begin{ex}
% ID: L11C2B15-03-TLN
% Mức: VD
Trong dạng “Lắp đặt và sử dụng ống cộng hưởng”, Các khoảng cách cộng hưởng đo được là $33,8\,\text{cm}$; $34,0\,\text{cm}$; $34,2\,\text{cm}$. Tính giá trị trung bình theo cm.
\shortans{34}
\loigiai{
$\bar d=(33,8+34,0+34,2)/3=34,0$ cm.
}
\end{ex}

% ===== Câu 12 =====
% ID: L11C2B15-03-TN
% Mức: TH
\begin{ex}
% ID: L11C2B15-03-TN
% Mức: TH
Trong dạng “Lắp đặt và sử dụng ống cộng hưởng”, Khoảng cách giữa hai vị trí cộng hưởng liên tiếp trong ống một đầu kín xấp xỉ
\choice
{\True $\lambda/2$}
{$\lambda$}
{$\lambda/4$}
{$2\lambda$}
\loigiai{
Đáp án A. Hai trạng thái cộng hưởng liên tiếp chênh nhau nửa bước sóng.
}
\end{ex}

% ===== Câu 13 =====
% ID: L11C2B15-04-DS
% Mức: VD
\begin{ex}
% ID: L11C2B15-04-DS
% Mức: VD
Xét các phát biểu sau trong dạng “Lắp đặt và sử dụng ống cộng hưởng”.
\choiceTF
{\True Khi đọc thước, mắt phải đặt — phát biểu đúng là: vuông góc với vạch chia tại vị trí cần đọc.}
{Khi đọc thước, mắt phải đặt — phát biểu càng xa thước càng tốt là đúng.}
{\True Nếu $f=500$ Hz và $\lambda=0,68$ m thì tốc độ âm là — phát biểu đúng là: $340\,\text{m}$ /s.}
{Nếu $f=500$ Hz và $\lambda=0,68$ m thì tốc độ âm là — phát biểu $0,00136\,\text{m}$ /s là đúng.}
\loigiai{
\begin{itemchoice}
\item (Đúng) Khi đọc thước, mắt phải đặt — phát biểu đúng là: vuông góc với vạch chia tại vị trí cần đọc. (Vì): ặt mắt vuông góc giúp tránh sai số thị sai. b) Sai: phương án nêu không phù hợp. c) Đúng: $v=f\lambda=500\cdot0,68=340$ m/s. d) Sai: phương án nêu không phù hợp
\item (Sai) Khi đọc thước, mắt phải đặt — phát biểu càng xa thước càng tốt là đúng.
\item (Đúng) Nếu $f=500$ Hz và $\lambda=0,68$ m thì tốc độ âm là — phát biểu đúng là: $340\,\text{m}$ /s.
\item (Sai) Nếu $f=500$ Hz và $\lambda=0,68$ m thì tốc độ âm là — phát biểu $0,00136\,\text{m}$ /s là đúng.
\end{itemchoice}
}
\end{ex}

% ===== Câu 14 =====
% ID: L11C2B15-04-TL
% Mức: VD
\begin{ex}
% ID: L11C2B15-04-TL
% Mức: VD
Trong dạng “Lắp đặt và sử dụng ống cộng hưởng”, Nêu hai biện pháp giảm sai số của phép đo.
\choiceTF

\loigiai{
Đặt mắt vuông góc khi đọc thước; xác định nhiều cặp cộng hưởng, đo lặp lại và lấy trung bình; giữ tần số, nhiệt độ ổn định.
}
\end{ex}

% ===== Câu 15 =====
% ID: L11C2B15-04-TLN
% Mức: VD
\begin{ex}
% ID: L11C2B15-04-TLN
% Mức: VD
Trong dạng “Lắp đặt và sử dụng ống cộng hưởng”, Với $\bar d=34$ cm và $f=500$ Hz, tính tốc độ âm theo m/s.
\shortans{340}
\loigiai{
$\lambda=2d=0,68$ m; $v=500\cdot0,68=340$ m/s.
}
\end{ex}

% ===== Câu 16 =====
% ID: L11C2B15-04-TN
% Mức: TH
\begin{ex}
% ID: L11C2B15-04-TN
% Mức: TH
Trong dạng “Lắp đặt và sử dụng ống cộng hưởng”, Biết khoảng cách hai vị trí cộng hưởng liên tiếp là d, bước sóng bằng
\choice
{\True $2d$}
{$d/2$}
{$d$}
{$4d$}
\loigiai{
Đáp án A. Vì $d=\lambda/2$ nên $\lambda=2d$.
}
\end{ex}

% ===== Câu 17 =====
% ID: L11C2B15-05-DS
% Mức: VDC
\begin{ex}
% ID: L11C2B15-05-DS
% Mức: VDC
Xét các phát biểu sau trong dạng “Lắp đặt và sử dụng ống cộng hưởng”.
\choiceTF
{\True Trong thí nghiệm, đại lượng thường được đo trực tiếp là — phát biểu đúng là: chiều dài cột khí tại cộng hưởng.}
{Trong thí nghiệm, đại lượng thường được đo trực tiếp là — phát biểu tốc độ âm là đúng.}
{\True Khi nhiệt độ không khí tăng, tốc độ truyền âm trong không khí nhìn chung — phát biểu đúng là: tăng.}
{Khi nhiệt độ không khí tăng, tốc độ truyền âm trong không khí nhìn chung — phát biểu không đổi là đúng.}
\loigiai{
\begin{itemchoice}
\item (Đúng) Trong thí nghiệm, đại lượng thường được đo trực tiếp là — phát biểu đúng là: chiều dài cột khí tại cộng hưởng. (Vì): Ta đọc trực tiếp chiều dài cột khí; bước sóng và tốc độ được tính gián tiếp. b) Sai: phương án nêu không phù hợp. c) Đúng: Tốc độ âm trong khí tăng khi nhiệt độ tăng. d) Sai: phương án nêu không phù hợp
\item (Sai) Trong thí nghiệm, đại lượng thường được đo trực tiếp là — phát biểu tốc độ âm là đúng.
\item (Đúng) Khi nhiệt độ không khí tăng, tốc độ truyền âm trong không khí nhìn chung — phát biểu đúng là: tăng.
\item (Sai) Khi nhiệt độ không khí tăng, tốc độ truyền âm trong không khí nhìn chung — phát biểu không đổi là đúng.
\end{itemchoice}
}
\end{ex}

% ===== Câu 18 =====
% ID: L11C2B15-05-TL
% Mức: VD
\begin{ex}
% ID: L11C2B15-05-TL
% Mức: VD
Trong dạng “Lắp đặt và sử dụng ống cộng hưởng”, Phân biệt đại lượng đo trực tiếp và gián tiếp trong thí nghiệm.
\choiceTF

\loigiai{
Chiều dài cột khí và tần số hiển thị là số liệu trực tiếp; bước sóng và tốc độ âm được suy ra bằng công thức nên là đại lượng gián tiếp.
}
\end{ex}

% ===== Câu 19 =====
% ID: L11C2B15-05-TLN
% Mức: VD
\begin{ex}
% ID: L11C2B15-05-TLN
% Mức: VD
Trong dạng “Lắp đặt và sử dụng ống cộng hưởng”, Giá trị đo là $338\,\text{m}$ /s, giá trị tham chiếu $340\,\text{m}$ /s. Tính độ lệch tuyệt đối theo m/s.
\shortans{2}
\loigiai{
$\Delta v=|338-340|=2$ m/s.
}
\end{ex}

% ===== Câu 20 =====
% ID: L11C2B15-05-TN
% Mức: TH
\begin{ex}
% ID: L11C2B15-05-TN
% Mức: TH
Trong dạng “Lắp đặt và sử dụng ống cộng hưởng”, Công thức tính tốc độ truyền âm là
\choice
{\True $v=f\lambda$}
{$v=f/\lambda$}
{$v=\lambda/f$}
{$v=f+\lambda$}
\loigiai{
Đáp án A. Tốc độ truyền sóng thỏa $v=f\lambda$.
}
\end{ex}

% ===== Câu 21 =====
% ID: L11C2B15-06-DS
% Mức: TH
\dangbt{Nhận biết mục đích và dụng cụ thí nghiệm đo tốc độ truyền âm}
\begin{ex}
% ID: L11C2B15-06-DS
% Mức: TH
Xét các phát biểu sau trong dạng “Nhận biết mục đích và dụng cụ thí nghiệm đo tốc độ truyền âm”.
\choiceTF
{\True Mục đích chính của thí nghiệm dùng ống cộng hưởng là xác định — phát biểu đúng là: tốc độ truyền âm trong không khí.}
{Mục đích chính của thí nghiệm dùng ống cộng hưởng là xác định — phát biểu cường độ dòng điện là đúng.}
{\True Dụng cụ phát âm có tần số xác định trong thí nghiệm thường là — phát biểu đúng là: âm thoa hoặc loa nối máy phát âm tần.}
{Dụng cụ phát âm có tần số xác định trong thí nghiệm thường là — phát biểu lực kế là đúng.}
\loigiai{
\begin{itemchoice}
\item (Đúng) Mục đích chính của thí nghiệm dùng ống cộng hưởng là xác định — phát biểu đúng là: tốc độ truyền âm trong không khí. (Vì): Từ tần số đã biết và bước sóng đo qua cộng hưởng, ta tính tốc độ truyền âm. b) Sai: phương án nêu không phù hợp. c) Đúng: Âm thoa hoặc loa tạo âm có tần số biết trước, cần cho phép đo. d) Sai: phương án nêu không phù hợp
\item (Sai) Mục đích chính của thí nghiệm dùng ống cộng hưởng là xác định — phát biểu cường độ dòng điện là đúng.
\item (Đúng) Dụng cụ phát âm có tần số xác định trong thí nghiệm thường là — phát biểu đúng là: âm thoa hoặc loa nối máy phát âm tần.
\item (Sai) Dụng cụ phát âm có tần số xác định trong thí nghiệm thường là — phát biểu lực kế là đúng.
\end{itemchoice}
}
\end{ex}

% ===== Câu 22 =====
% ID: L11C2B15-06-TL
% Mức: VDC
\dangbt{Lắp đặt và sử dụng ống cộng hưởng}
\begin{ex}
% ID: L11C2B15-06-TL
% Mức: VDC
Trong dạng “Lắp đặt và sử dụng ống cộng hưởng”, Một nhóm đo được $342\,\text{m}$ /s trong khi giá trị tham chiếu là $340\,\text{m}$ /s. Nhận xét kết quả.
\choiceTF

\loigiai{
Độ lệch tuyệt đối là $2\,\text{m}$ /s, sai lệch tương đối khoảng $2/340\cdot100\%=0,59\%$; kết quả khá phù hợp trong điều kiện thí nghiệm phổ thông.
}
\end{ex}

% ===== Câu 23 =====
% ID: L11C2B15-06-TLN
% Mức: VDC
\begin{ex}
% ID: L11C2B15-06-TLN
% Mức: VDC
Trong dạng “Lắp đặt và sử dụng ống cộng hưởng”, Đo được $d=0,30$ m với âm $600\,\text{Hz}$. Tính tốc độ âm theo m/s.
\shortans{360}
\loigiai{
$\lambda=2d=0,60$ m; $v=f\lambda=600\cdot0,60=360$ m/s.
}
\end{ex}

% ===== Câu 24 =====
% ID: L11C2B15-06-TN
% Mức: VD
\begin{ex}
% ID: L11C2B15-06-TN
% Mức: VD
Trong dạng “Lắp đặt và sử dụng ống cộng hưởng”, Để giảm sai số khi xác định khoảng cách cộng hưởng nên
\choice
{\True đo nhiều cặp vị trí và lấy trung bình}
{chỉ đo một lần}
{thay đổi tần số liên tục trong mỗi lần đo}
{đọc thước từ góc xiên}
\loigiai{
Đáp án A. Đo lặp lại và lấy trung bình làm giảm ảnh hưởng sai số ngẫu nhiên.
}
\end{ex}

% ===== Câu 25 =====
% ID: L11C2B15-07-DS
% Mức: TH
\dangbt{Nhận biết mục đích và dụng cụ thí nghiệm đo tốc độ truyền âm}
\begin{ex}
% ID: L11C2B15-07-DS
% Mức: TH
Xét các phát biểu sau trong dạng “Nhận biết mục đích và dụng cụ thí nghiệm đo tốc độ truyền âm”.
\choiceTF
{\True Khoảng cách giữa hai vị trí cộng hưởng liên tiếp trong ống một đầu kín xấp xỉ — phát biểu đúng là: $\lambda/2$.}
{Khoảng cách giữa hai vị trí cộng hưởng liên tiếp trong ống một đầu kín xấp xỉ — phát biểu $\lambda$ là đúng.}
{\True Biết khoảng cách hai vị trí cộng hưởng liên tiếp là d, bước sóng bằng — phát biểu đúng là: $2d$.}
{Biết khoảng cách hai vị trí cộng hưởng liên tiếp là d, bước sóng bằng — phát biểu $d$ là đúng.}
\loigiai{
\begin{itemchoice}
\item (Đúng) Khoảng cách giữa hai vị trí cộng hưởng liên tiếp trong ống một đầu kín xấp xỉ — phát biểu đúng là: $\lambda/2$. (Vì): Hai trạng thái cộng hưởng liên tiếp chênh nhau nửa bước sóng. b) Sai: phương án nêu không phù hợp. c) Đúng: Vì $d=\lambda/2$ nên $\lambda=2d$. d) Sai: phương án nêu không phù hợp
\item (Sai) Khoảng cách giữa hai vị trí cộng hưởng liên tiếp trong ống một đầu kín xấp xỉ — phát biểu $\lambda$ là đúng.
\item (Đúng) Biết khoảng cách hai vị trí cộng hưởng liên tiếp là d, bước sóng bằng — phát biểu đúng là: $2d$.
\item (Sai) Biết khoảng cách hai vị trí cộng hưởng liên tiếp là d, bước sóng bằng — phát biểu $d$ là đúng.
\end{itemchoice}
}
\end{ex}

% ===== Câu 26 =====
% ID: L11C2B15-07-TL
% Mức: TH
\begin{ex}
% ID: L11C2B15-07-TL
% Mức: TH
Trong dạng “Nhận biết mục đích và dụng cụ thí nghiệm đo tốc độ truyền âm”, Nêu các bước chính đo tốc độ truyền âm bằng ống cộng hưởng.
\choiceTF

\loigiai{
Phát âm có tần số f đã biết; thay đổi chiều dài cột khí; ghi hai vị trí cộng hưởng liên tiếp; tính $\lambda=2|L_2-L_1|$ rồi tính $v=f\lambda$.
}
\end{ex}

% ===== Câu 27 =====
% ID: L11C2B15-07-TLN
% Mức: TH
\begin{ex}
% ID: L11C2B15-07-TLN
% Mức: TH
Trong dạng “Nhận biết mục đích và dụng cụ thí nghiệm đo tốc độ truyền âm”, Hai vị trí cộng hưởng liên tiếp cách nhau $34\,\text{cm}$. Tính bước sóng theo cm.
\shortans{68}
\loigiai{
$\lambda=2d=2\cdot34=68$ cm.
}
\end{ex}

% ===== Câu 28 =====
% ID: L11C2B15-07-TN
% Mức: VD
\dangbt{Lắp đặt và sử dụng ống cộng hưởng}
\begin{ex}
% ID: L11C2B15-07-TN
% Mức: VD
Trong dạng “Lắp đặt và sử dụng ống cộng hưởng”, Khi đọc thước, mắt phải đặt
\choice
{\True vuông góc với vạch chia tại vị trí cần đọc}
{càng xa thước càng tốt}
{xiên 45 độ}
{ở bất kì vị trí nào}
\loigiai{
Đáp án A. Đặt mắt vuông góc giúp tránh sai số thị sai.
}
\end{ex}

% ===== Câu 29 =====
% ID: L11C2B15-08-DS
% Mức: VD
\dangbt{Nhận biết mục đích và dụng cụ thí nghiệm đo tốc độ truyền âm}
\begin{ex}
% ID: L11C2B15-08-DS
% Mức: VD
Xét các phát biểu sau trong dạng “Nhận biết mục đích và dụng cụ thí nghiệm đo tốc độ truyền âm”.
\choiceTF
{\True Công thức tính tốc độ truyền âm là — phát biểu đúng là: $v=f\lambda$.}
{Công thức tính tốc độ truyền âm là — phát biểu $v=f/\lambda$ là đúng.}
{\True Để giảm sai số khi xác định khoảng cách cộng hưởng nên — phát biểu đúng là: đo nhiều cặp vị trí và lấy trung bình.}
{Để giảm sai số khi xác định khoảng cách cộng hưởng nên — phát biểu thay đổi tần số liên tục trong mỗi lần đo là đúng.}
\loigiai{
\begin{itemchoice}
\item (Đúng) Công thức tính tốc độ truyền âm là — phát biểu đúng là: $v=f\lambda$. (Vì): Tốc độ truyền sóng thỏa $v=f\lambda$. b) Sai: phương án nêu không phù hợp. c) Đúng: Đo lặp lại và lấy trung bình làm giảm ảnh hưởng sai số ngẫu nhiên. d) Sai: phương án nêu không phù hợp
\item (Sai) Công thức tính tốc độ truyền âm là — phát biểu $v=f/\lambda$ là đúng.
\item (Đúng) Để giảm sai số khi xác định khoảng cách cộng hưởng nên — phát biểu đúng là: đo nhiều cặp vị trí và lấy trung bình.
\item (Sai) Để giảm sai số khi xác định khoảng cách cộng hưởng nên — phát biểu thay đổi tần số liên tục trong mỗi lần đo là đúng.
\end{itemchoice}
}
\end{ex}

% ===== Câu 30 =====
% ID: L11C2B15-08-TL
% Mức: TH
\begin{ex}
% ID: L11C2B15-08-TL
% Mức: TH
Trong dạng “Nhận biết mục đích và dụng cụ thí nghiệm đo tốc độ truyền âm”, Giải thích vì sao dùng hai vị trí cộng hưởng liên tiếp giúp xác định bước sóng.
\choiceTF

\loigiai{
Trong ống một đầu kín, các chiều dài cộng hưởng kế tiếp chênh nhau $\lambda/2$, nên đo hiệu hai vị trí loại được phần hiệu chỉnh đầu ống và cho $\lambda=2\Delta L$.
}
\end{ex}

% ===== Câu 31 =====
% ID: L11C2B15-08-TLN
% Mức: TH
\begin{ex}
% ID: L11C2B15-08-TLN
% Mức: TH
Trong dạng “Nhận biết mục đích và dụng cụ thí nghiệm đo tốc độ truyền âm”, Âm có tần số $500\,\text{Hz}$, bước sóng $0,68\,\text{m}$. Tính tốc độ theo m/s.
\shortans{340}
\loigiai{
$v=f\lambda=500\cdot0,68=340$ m/s.
}
\end{ex}

% ===== Câu 32 =====
% ID: L11C2B15-08-TN
% Mức: VD
\dangbt{Lắp đặt và sử dụng ống cộng hưởng}
\begin{ex}
% ID: L11C2B15-08-TN
% Mức: VD
Trong dạng “Lắp đặt và sử dụng ống cộng hưởng”, Nếu $f=500$ Hz và $\lambda=0,68$ m thì tốc độ âm là
\choice
{\True $340\,\text{m}$ /s}
{$735\,\text{m}$ /s}
{$0,00136\,\text{m}$ /s}
{$500,68\,\text{m}$ /s}
\loigiai{
Đáp án A. $v=f\lambda=500\cdot0,68=340$ m/s.
}
\end{ex}

% ===== Câu 33 =====
% ID: L11C2B15-09-DS
% Mức: VD
\dangbt{Nhận biết mục đích và dụng cụ thí nghiệm đo tốc độ truyền âm}
\begin{ex}
% ID: L11C2B15-09-DS
% Mức: VD
Xét các phát biểu sau trong dạng “Nhận biết mục đích và dụng cụ thí nghiệm đo tốc độ truyền âm”.
\choiceTF
{\True Khi đọc thước, mắt phải đặt — phát biểu đúng là: vuông góc với vạch chia tại vị trí cần đọc.}
{Khi đọc thước, mắt phải đặt — phát biểu càng xa thước càng tốt là đúng.}
{\True Nếu $f=500$ Hz và $\lambda=0,68$ m thì tốc độ âm là — phát biểu đúng là: $340\,\text{m}$ /s.}
{Nếu $f=500$ Hz và $\lambda=0,68$ m thì tốc độ âm là — phát biểu $0,00136\,\text{m}$ /s là đúng.}
\loigiai{
\begin{itemchoice}
\item (Đúng) Khi đọc thước, mắt phải đặt — phát biểu đúng là: vuông góc với vạch chia tại vị trí cần đọc. (Vì): ặt mắt vuông góc giúp tránh sai số thị sai. b) Sai: phương án nêu không phù hợp. c) Đúng: $v=f\lambda=500\cdot0,68=340$ m/s. d) Sai: phương án nêu không phù hợp
\item (Sai) Khi đọc thước, mắt phải đặt — phát biểu càng xa thước càng tốt là đúng.
\item (Đúng) Nếu $f=500$ Hz và $\lambda=0,68$ m thì tốc độ âm là — phát biểu đúng là: $340\,\text{m}$ /s.
\item (Sai) Nếu $f=500$ Hz và $\lambda=0,68$ m thì tốc độ âm là — phát biểu $0,00136\,\text{m}$ /s là đúng.
\end{itemchoice}
}
\end{ex}

% ===== Câu 34 =====
% ID: L11C2B15-09-TL
% Mức: VD
\begin{ex}
% ID: L11C2B15-09-TL
% Mức: VD
Trong dạng “Nhận biết mục đích và dụng cụ thí nghiệm đo tốc độ truyền âm”, Với $f=400$ Hz, hai vị trí cộng hưởng là $18\,\text{cm}$ và $60\,\text{cm}$. Tính tốc độ âm.
\choiceTF

\loigiai{
$d=60-18=42\,\mathrm{cm}=0,42$ m; $\lambda=0,84$ m; $v=400\cdot0,84=336$ m/s.
}
\end{ex}

% ===== Câu 35 =====
% ID: L11C2B15-09-TLN
% Mức: VD
\begin{ex}
% ID: L11C2B15-09-TLN
% Mức: VD
Trong dạng “Nhận biết mục đích và dụng cụ thí nghiệm đo tốc độ truyền âm”, Các khoảng cách cộng hưởng đo được là $33,8\,\text{cm}$; $34,0\,\text{cm}$; $34,2\,\text{cm}$. Tính giá trị trung bình theo cm.
\shortans{34}
\loigiai{
$\bar d=(33,8+34,0+34,2)/3=34,0$ cm.
}
\end{ex}

% ===== Câu 36 =====
% ID: L11C2B15-09-TN
% Mức: VDC
\dangbt{Lắp đặt và sử dụng ống cộng hưởng}
\begin{ex}
% ID: L11C2B15-09-TN
% Mức: VDC
Trong dạng “Lắp đặt và sử dụng ống cộng hưởng”, Trong thí nghiệm, đại lượng thường được đo trực tiếp là
\choice
{\True chiều dài cột khí tại cộng hưởng}
{tốc độ âm}
{bước sóng chính xác tuyệt đối}
{pha ban đầu của nguồn}
\loigiai{
Đáp án A. Ta đọc trực tiếp chiều dài cột khí; bước sóng và tốc độ được tính gián tiếp.
}
\end{ex}

% ===== Câu 37 =====
% ID: L11C2B15-10-DS
% Mức: VDC
\dangbt{Nhận biết mục đích và dụng cụ thí nghiệm đo tốc độ truyền âm}
\begin{ex}
% ID: L11C2B15-10-DS
% Mức: VDC
Xét các phát biểu sau trong dạng “Nhận biết mục đích và dụng cụ thí nghiệm đo tốc độ truyền âm”.
\choiceTF
{\True Trong thí nghiệm, đại lượng thường được đo trực tiếp là — phát biểu đúng là: chiều dài cột khí tại cộng hưởng.}
{Trong thí nghiệm, đại lượng thường được đo trực tiếp là — phát biểu tốc độ âm là đúng.}
{\True Khi nhiệt độ không khí tăng, tốc độ truyền âm trong không khí nhìn chung — phát biểu đúng là: tăng.}
{Khi nhiệt độ không khí tăng, tốc độ truyền âm trong không khí nhìn chung — phát biểu không đổi là đúng.}
\loigiai{
\begin{itemchoice}
\item (Đúng) Trong thí nghiệm, đại lượng thường được đo trực tiếp là — phát biểu đúng là: chiều dài cột khí tại cộng hưởng. (Vì): Ta đọc trực tiếp chiều dài cột khí; bước sóng và tốc độ được tính gián tiếp. b) Sai: phương án nêu không phù hợp. c) Đúng: Tốc độ âm trong khí tăng khi nhiệt độ tăng. d) Sai: phương án nêu không phù hợp
\item (Sai) Trong thí nghiệm, đại lượng thường được đo trực tiếp là — phát biểu tốc độ âm là đúng.
\item (Đúng) Khi nhiệt độ không khí tăng, tốc độ truyền âm trong không khí nhìn chung — phát biểu đúng là: tăng.
\item (Sai) Khi nhiệt độ không khí tăng, tốc độ truyền âm trong không khí nhìn chung — phát biểu không đổi là đúng.
\end{itemchoice}
}
\end{ex}

% ===== Câu 38 =====
% ID: L11C2B15-10-TL
% Mức: VD
\begin{ex}
% ID: L11C2B15-10-TL
% Mức: VD
Trong dạng “Nhận biết mục đích và dụng cụ thí nghiệm đo tốc độ truyền âm”, Nêu hai biện pháp giảm sai số của phép đo.
\choiceTF

\loigiai{
Đặt mắt vuông góc khi đọc thước; xác định nhiều cặp cộng hưởng, đo lặp lại và lấy trung bình; giữ tần số, nhiệt độ ổn định.
}
\end{ex}

% ===== Câu 39 =====
% ID: L11C2B15-10-TLN
% Mức: VD
\begin{ex}
% ID: L11C2B15-10-TLN
% Mức: VD
Trong dạng “Nhận biết mục đích và dụng cụ thí nghiệm đo tốc độ truyền âm”, Với $\bar d=34$ cm và $f=500$ Hz, tính tốc độ âm theo m/s.
\shortans{340}
\loigiai{
$\lambda=2d=0,68$ m; $v=500\cdot0,68=340$ m/s.
}
\end{ex}

% ===== Câu 40 =====
% ID: L11C2B15-10-TN
% Mức: VDC
\dangbt{Lắp đặt và sử dụng ống cộng hưởng}
\begin{ex}
% ID: L11C2B15-10-TN
% Mức: VDC
Trong dạng “Lắp đặt và sử dụng ống cộng hưởng”, Khi nhiệt độ không khí tăng, tốc độ truyền âm trong không khí nhìn chung
\choice
{\True tăng}
{giảm}
{không đổi}
{bằng 0}
\loigiai{
Đáp án A. Tốc độ âm trong khí tăng khi nhiệt độ tăng.
}
\end{ex}

% ===== Câu 41 =====
% ID: L11C2B15-11-DS
% Mức: TH
\dangbt{Tính tốc độ truyền âm từ số liệu thí nghiệm}
\begin{ex}
% ID: L11C2B15-11-DS
% Mức: TH
Xét các phát biểu sau trong dạng “Tính tốc độ truyền âm từ số liệu thí nghiệm”.
\choiceTF
{\True Mục đích chính của thí nghiệm dùng ống cộng hưởng là xác định — phát biểu đúng là: tốc độ truyền âm trong không khí.}
{Mục đích chính của thí nghiệm dùng ống cộng hưởng là xác định — phát biểu cường độ dòng điện là đúng.}
{\True Dụng cụ phát âm có tần số xác định trong thí nghiệm thường là — phát biểu đúng là: âm thoa hoặc loa nối máy phát âm tần.}
{Dụng cụ phát âm có tần số xác định trong thí nghiệm thường là — phát biểu lực kế là đúng.}
\loigiai{
\begin{itemchoice}
\item (Đúng) Mục đích chính của thí nghiệm dùng ống cộng hưởng là xác định — phát biểu đúng là: tốc độ truyền âm trong không khí. (Vì): Từ tần số đã biết và bước sóng đo qua cộng hưởng, ta tính tốc độ truyền âm. b) Sai: phương án nêu không phù hợp. c) Đúng: Âm thoa hoặc loa tạo âm có tần số biết trước, cần cho phép đo. d) Sai: phương án nêu không phù hợp
\item (Sai) Mục đích chính của thí nghiệm dùng ống cộng hưởng là xác định — phát biểu cường độ dòng điện là đúng.
\item (Đúng) Dụng cụ phát âm có tần số xác định trong thí nghiệm thường là — phát biểu đúng là: âm thoa hoặc loa nối máy phát âm tần.
\item (Sai) Dụng cụ phát âm có tần số xác định trong thí nghiệm thường là — phát biểu lực kế là đúng.
\end{itemchoice}
}
\end{ex}

% ===== Câu 42 =====
% ID: L11C2B15-11-TL
% Mức: VD
\dangbt{Nhận biết mục đích và dụng cụ thí nghiệm đo tốc độ truyền âm}
\begin{ex}
% ID: L11C2B15-11-TL
% Mức: VD
Trong dạng “Nhận biết mục đích và dụng cụ thí nghiệm đo tốc độ truyền âm”, Phân biệt đại lượng đo trực tiếp và gián tiếp trong thí nghiệm.
\choiceTF

\loigiai{
Chiều dài cột khí và tần số hiển thị là số liệu trực tiếp; bước sóng và tốc độ âm được suy ra bằng công thức nên là đại lượng gián tiếp.
}
\end{ex}

% ===== Câu 43 =====
% ID: L11C2B15-11-TLN
% Mức: VD
\begin{ex}
% ID: L11C2B15-11-TLN
% Mức: VD
Trong dạng “Nhận biết mục đích và dụng cụ thí nghiệm đo tốc độ truyền âm”, Giá trị đo là $338\,\text{m}$ /s, giá trị tham chiếu $340\,\text{m}$ /s. Tính độ lệch tuyệt đối theo m/s.
\shortans{2}
\loigiai{
$\Delta v=|338-340|=2$ m/s.
}
\end{ex}

% ===== Câu 44 =====
% ID: L11C2B15-11-TN
% Mức: NB
\begin{ex}
% ID: L11C2B15-11-TN
% Mức: NB
Trong dạng “Nhận biết mục đích và dụng cụ thí nghiệm đo tốc độ truyền âm”, Mục đích chính của thí nghiệm dùng ống cộng hưởng là xác định
\choice
{\True tốc độ truyền âm trong không khí}
{cường độ dòng điện}
{khối lượng riêng của nước}
{nhiệt dung riêng của không khí}
\loigiai{
Đáp án A. Từ tần số đã biết và bước sóng đo qua cộng hưởng, ta tính tốc độ truyền âm.
}
\end{ex}

% ===== Câu 45 =====
% ID: L11C2B15-12-DS
% Mức: TH
\dangbt{Tính tốc độ truyền âm từ số liệu thí nghiệm}
\begin{ex}
% ID: L11C2B15-12-DS
% Mức: TH
Xét các phát biểu sau trong dạng “Tính tốc độ truyền âm từ số liệu thí nghiệm”.
\choiceTF
{\True Khoảng cách giữa hai vị trí cộng hưởng liên tiếp trong ống một đầu kín xấp xỉ — phát biểu đúng là: $\lambda/2$.}
{Khoảng cách giữa hai vị trí cộng hưởng liên tiếp trong ống một đầu kín xấp xỉ — phát biểu $\lambda$ là đúng.}
{\True Biết khoảng cách hai vị trí cộng hưởng liên tiếp là d, bước sóng bằng — phát biểu đúng là: $2d$.}
{Biết khoảng cách hai vị trí cộng hưởng liên tiếp là d, bước sóng bằng — phát biểu $d$ là đúng.}
\loigiai{
\begin{itemchoice}
\item (Đúng) Khoảng cách giữa hai vị trí cộng hưởng liên tiếp trong ống một đầu kín xấp xỉ — phát biểu đúng là: $\lambda/2$. (Vì): Hai trạng thái cộng hưởng liên tiếp chênh nhau nửa bước sóng. b) Sai: phương án nêu không phù hợp. c) Đúng: Vì $d=\lambda/2$ nên $\lambda=2d$. d) Sai: phương án nêu không phù hợp
\item (Sai) Khoảng cách giữa hai vị trí cộng hưởng liên tiếp trong ống một đầu kín xấp xỉ — phát biểu $\lambda$ là đúng.
\item (Đúng) Biết khoảng cách hai vị trí cộng hưởng liên tiếp là d, bước sóng bằng — phát biểu đúng là: $2d$.
\item (Sai) Biết khoảng cách hai vị trí cộng hưởng liên tiếp là d, bước sóng bằng — phát biểu $d$ là đúng.
\end{itemchoice}
}
\end{ex}

% ===== Câu 46 =====
% ID: L11C2B15-12-TL
% Mức: VDC
\dangbt{Nhận biết mục đích và dụng cụ thí nghiệm đo tốc độ truyền âm}
\begin{ex}
% ID: L11C2B15-12-TL
% Mức: VDC
Trong dạng “Nhận biết mục đích và dụng cụ thí nghiệm đo tốc độ truyền âm”, Một nhóm đo được $342\,\text{m}$ /s trong khi giá trị tham chiếu là $340\,\text{m}$ /s. Nhận xét kết quả.
\choiceTF

\loigiai{
Độ lệch tuyệt đối là $2\,\text{m}$ /s, sai lệch tương đối khoảng $2/340\cdot100\%=0,59\%$; kết quả khá phù hợp trong điều kiện thí nghiệm phổ thông.
}
\end{ex}

% ===== Câu 47 =====
% ID: L11C2B15-12-TLN
% Mức: VDC
\begin{ex}
% ID: L11C2B15-12-TLN
% Mức: VDC
Trong dạng “Nhận biết mục đích và dụng cụ thí nghiệm đo tốc độ truyền âm”, Đo được $d=0,30$ m với âm $600\,\text{Hz}$. Tính tốc độ âm theo m/s.
\shortans{360}
\loigiai{
$\lambda=2d=0,60$ m; $v=f\lambda=600\cdot0,60=360$ m/s.
}
\end{ex}

% ===== Câu 48 =====
% ID: L11C2B15-12-TN
% Mức: NB
\begin{ex}
% ID: L11C2B15-12-TN
% Mức: NB
Trong dạng “Nhận biết mục đích và dụng cụ thí nghiệm đo tốc độ truyền âm”, Dụng cụ phát âm có tần số xác định trong thí nghiệm thường là
\choice
{\True âm thoa hoặc loa nối máy phát âm tần}
{nhiệt kế}
{lực kế}
{cân điện tử}
\loigiai{
Đáp án A. Âm thoa hoặc loa tạo âm có tần số biết trước, cần cho phép đo.
}
\end{ex}

% ===== Câu 49 =====
% ID: L11C2B15-13-DS
% Mức: VD
\dangbt{Tính tốc độ truyền âm từ số liệu thí nghiệm}
\begin{ex}
% ID: L11C2B15-13-DS
% Mức: VD
Xét các phát biểu sau trong dạng “Tính tốc độ truyền âm từ số liệu thí nghiệm”.
\choiceTF
{\True Công thức tính tốc độ truyền âm là — phát biểu đúng là: $v=f\lambda$.}
{Công thức tính tốc độ truyền âm là — phát biểu $v=f/\lambda$ là đúng.}
{\True Để giảm sai số khi xác định khoảng cách cộng hưởng nên — phát biểu đúng là: đo nhiều cặp vị trí và lấy trung bình.}
{Để giảm sai số khi xác định khoảng cách cộng hưởng nên — phát biểu thay đổi tần số liên tục trong mỗi lần đo là đúng.}
\loigiai{
\begin{itemchoice}
\item (Đúng) Công thức tính tốc độ truyền âm là — phát biểu đúng là: $v=f\lambda$. (Vì): Tốc độ truyền sóng thỏa $v=f\lambda$. b) Sai: phương án nêu không phù hợp. c) Đúng: Đo lặp lại và lấy trung bình làm giảm ảnh hưởng sai số ngẫu nhiên. d) Sai: phương án nêu không phù hợp
\item (Sai) Công thức tính tốc độ truyền âm là — phát biểu $v=f/\lambda$ là đúng.
\item (Đúng) Để giảm sai số khi xác định khoảng cách cộng hưởng nên — phát biểu đúng là: đo nhiều cặp vị trí và lấy trung bình.
\item (Sai) Để giảm sai số khi xác định khoảng cách cộng hưởng nên — phát biểu thay đổi tần số liên tục trong mỗi lần đo là đúng.
\end{itemchoice}
}
\end{ex}

% ===== Câu 50 =====
% ID: L11C2B15-13-TL
% Mức: TH
\begin{ex}
% ID: L11C2B15-13-TL
% Mức: TH
Trong dạng “Tính tốc độ truyền âm từ số liệu thí nghiệm”, Nêu các bước chính đo tốc độ truyền âm bằng ống cộng hưởng.
\choiceTF

\loigiai{
Phát âm có tần số f đã biết; thay đổi chiều dài cột khí; ghi hai vị trí cộng hưởng liên tiếp; tính $\lambda=2|L_2-L_1|$ rồi tính $v=f\lambda$.
}
\end{ex}

% ===== Câu 51 =====
% ID: L11C2B15-13-TLN
% Mức: TH
\begin{ex}
% ID: L11C2B15-13-TLN
% Mức: TH
Trong dạng “Tính tốc độ truyền âm từ số liệu thí nghiệm”, Hai vị trí cộng hưởng liên tiếp cách nhau $34\,\text{cm}$. Tính bước sóng theo cm.
\shortans{68}
\loigiai{
$\lambda=2d=2\cdot34=68$ cm.
}
\end{ex}

% ===== Câu 52 =====
% ID: L11C2B15-13-TN
% Mức: TH
\dangbt{Nhận biết mục đích và dụng cụ thí nghiệm đo tốc độ truyền âm}
\begin{ex}
% ID: L11C2B15-13-TN
% Mức: TH
Trong dạng “Nhận biết mục đích và dụng cụ thí nghiệm đo tốc độ truyền âm”, Khoảng cách giữa hai vị trí cộng hưởng liên tiếp trong ống một đầu kín xấp xỉ
\choice
{\True $\lambda/2$}
{$\lambda$}
{$\lambda/4$}
{$2\lambda$}
\loigiai{
Đáp án A. Hai trạng thái cộng hưởng liên tiếp chênh nhau nửa bước sóng.
}
\end{ex}

% ===== Câu 53 =====
% ID: L11C2B15-14-DS
% Mức: VD
\dangbt{Tính tốc độ truyền âm từ số liệu thí nghiệm}
\begin{ex}
% ID: L11C2B15-14-DS
% Mức: VD
Xét các phát biểu sau trong dạng “Tính tốc độ truyền âm từ số liệu thí nghiệm”.
\choiceTF
{\True Khi đọc thước, mắt phải đặt — phát biểu đúng là: vuông góc với vạch chia tại vị trí cần đọc.}
{Khi đọc thước, mắt phải đặt — phát biểu càng xa thước càng tốt là đúng.}
{\True Nếu $f=500$ Hz và $\lambda=0,68$ m thì tốc độ âm là — phát biểu đúng là: $340\,\text{m}$ /s.}
{Nếu $f=500$ Hz và $\lambda=0,68$ m thì tốc độ âm là — phát biểu $0,00136\,\text{m}$ /s là đúng.}
\loigiai{
\begin{itemchoice}
\item (Đúng) Khi đọc thước, mắt phải đặt — phát biểu đúng là: vuông góc với vạch chia tại vị trí cần đọc. (Vì): ặt mắt vuông góc giúp tránh sai số thị sai. b) Sai: phương án nêu không phù hợp. c) Đúng: $v=f\lambda=500\cdot0,68=340$ m/s. d) Sai: phương án nêu không phù hợp
\item (Sai) Khi đọc thước, mắt phải đặt — phát biểu càng xa thước càng tốt là đúng.
\item (Đúng) Nếu $f=500$ Hz và $\lambda=0,68$ m thì tốc độ âm là — phát biểu đúng là: $340\,\text{m}$ /s.
\item (Sai) Nếu $f=500$ Hz và $\lambda=0,68$ m thì tốc độ âm là — phát biểu $0,00136\,\text{m}$ /s là đúng.
\end{itemchoice}
}
\end{ex}

% ===== Câu 54 =====
% ID: L11C2B15-14-TL
% Mức: TH
\begin{ex}
% ID: L11C2B15-14-TL
% Mức: TH
Trong dạng “Tính tốc độ truyền âm từ số liệu thí nghiệm”, Giải thích vì sao dùng hai vị trí cộng hưởng liên tiếp giúp xác định bước sóng.
\choiceTF

\loigiai{
Trong ống một đầu kín, các chiều dài cộng hưởng kế tiếp chênh nhau $\lambda/2$, nên đo hiệu hai vị trí loại được phần hiệu chỉnh đầu ống và cho $\lambda=2\Delta L$.
}
\end{ex}

% ===== Câu 55 =====
% ID: L11C2B15-14-TLN
% Mức: TH
\begin{ex}
% ID: L11C2B15-14-TLN
% Mức: TH
Trong dạng “Tính tốc độ truyền âm từ số liệu thí nghiệm”, Âm có tần số $500\,\text{Hz}$, bước sóng $0,68\,\text{m}$. Tính tốc độ theo m/s.
\shortans{340}
\loigiai{
$v=f\lambda=500\cdot0,68=340$ m/s.
}
\end{ex}

% ===== Câu 56 =====
% ID: L11C2B15-14-TN
% Mức: TH
\dangbt{Nhận biết mục đích và dụng cụ thí nghiệm đo tốc độ truyền âm}
\begin{ex}
% ID: L11C2B15-14-TN
% Mức: TH
Trong dạng “Nhận biết mục đích và dụng cụ thí nghiệm đo tốc độ truyền âm”, Biết khoảng cách hai vị trí cộng hưởng liên tiếp là d, bước sóng bằng
\choice
{\True $2d$}
{$d/2$}
{$d$}
{$4d$}
\loigiai{
Đáp án A. Vì $d=\lambda/2$ nên $\lambda=2d$.
}
\end{ex}

% ===== Câu 57 =====
% ID: L11C2B15-15-DS
% Mức: VDC
\dangbt{Tính tốc độ truyền âm từ số liệu thí nghiệm}
\begin{ex}
% ID: L11C2B15-15-DS
% Mức: VDC
Xét các phát biểu sau trong dạng “Tính tốc độ truyền âm từ số liệu thí nghiệm”.
\choiceTF
{\True Trong thí nghiệm, đại lượng thường được đo trực tiếp là — phát biểu đúng là: chiều dài cột khí tại cộng hưởng.}
{Trong thí nghiệm, đại lượng thường được đo trực tiếp là — phát biểu tốc độ âm là đúng.}
{\True Khi nhiệt độ không khí tăng, tốc độ truyền âm trong không khí nhìn chung — phát biểu đúng là: tăng.}
{Khi nhiệt độ không khí tăng, tốc độ truyền âm trong không khí nhìn chung — phát biểu không đổi là đúng.}
\loigiai{
\begin{itemchoice}
\item (Đúng) Trong thí nghiệm, đại lượng thường được đo trực tiếp là — phát biểu đúng là: chiều dài cột khí tại cộng hưởng. (Vì): Ta đọc trực tiếp chiều dài cột khí; bước sóng và tốc độ được tính gián tiếp. b) Sai: phương án nêu không phù hợp. c) Đúng: Tốc độ âm trong khí tăng khi nhiệt độ tăng. d) Sai: phương án nêu không phù hợp
\item (Sai) Trong thí nghiệm, đại lượng thường được đo trực tiếp là — phát biểu tốc độ âm là đúng.
\item (Đúng) Khi nhiệt độ không khí tăng, tốc độ truyền âm trong không khí nhìn chung — phát biểu đúng là: tăng.
\item (Sai) Khi nhiệt độ không khí tăng, tốc độ truyền âm trong không khí nhìn chung — phát biểu không đổi là đúng.
\end{itemchoice}
}
\end{ex}

% ===== Câu 58 =====
% ID: L11C2B15-15-TL
% Mức: VD
\begin{ex}
% ID: L11C2B15-15-TL
% Mức: VD
Trong dạng “Tính tốc độ truyền âm từ số liệu thí nghiệm”, Với $f=400$ Hz, hai vị trí cộng hưởng là $18\,\text{cm}$ và $60\,\text{cm}$. Tính tốc độ âm.
\choiceTF

\loigiai{
$d=60-18=42\,\mathrm{cm}=0,42$ m; $\lambda=0,84$ m; $v=400\cdot0,84=336$ m/s.
}
\end{ex}

% ===== Câu 59 =====
% ID: L11C2B15-15-TLN
% Mức: VD
\begin{ex}
% ID: L11C2B15-15-TLN
% Mức: VD
Trong dạng “Tính tốc độ truyền âm từ số liệu thí nghiệm”, Các khoảng cách cộng hưởng đo được là $33,8\,\text{cm}$; $34,0\,\text{cm}$; $34,2\,\text{cm}$. Tính giá trị trung bình theo cm.
\shortans{34}
\loigiai{
$\bar d=(33,8+34,0+34,2)/3=34,0$ cm.
}
\end{ex}

% ===== Câu 60 =====
% ID: L11C2B15-15-TN
% Mức: TH
\dangbt{Nhận biết mục đích và dụng cụ thí nghiệm đo tốc độ truyền âm}
\begin{ex}
% ID: L11C2B15-15-TN
% Mức: TH
Trong dạng “Nhận biết mục đích và dụng cụ thí nghiệm đo tốc độ truyền âm”, Công thức tính tốc độ truyền âm là
\choice
{\True $v=f\lambda$}
{$v=f/\lambda$}
{$v=\lambda/f$}
{$v=f+\lambda$}
\loigiai{
Đáp án A. Tốc độ truyền sóng thỏa $v=f\lambda$.
}
\end{ex}

% ===== Câu 61 =====
% ID: L11C2B15-16-DS
% Mức: TH
\dangbt{Xác định các vị trí cộng hưởng liên tiếp}
\begin{ex}
% ID: L11C2B15-16-DS
% Mức: TH
Xét các phát biểu sau trong dạng “Xác định các vị trí cộng hưởng liên tiếp”.
\choiceTF
{\True Mục đích chính của thí nghiệm dùng ống cộng hưởng là xác định — phát biểu đúng là: tốc độ truyền âm trong không khí.}
{Mục đích chính của thí nghiệm dùng ống cộng hưởng là xác định — phát biểu cường độ dòng điện là đúng.}
{\True Dụng cụ phát âm có tần số xác định trong thí nghiệm thường là — phát biểu đúng là: âm thoa hoặc loa nối máy phát âm tần.}
{Dụng cụ phát âm có tần số xác định trong thí nghiệm thường là — phát biểu lực kế là đúng.}
\loigiai{
\begin{itemchoice}
\item (Đúng) Mục đích chính của thí nghiệm dùng ống cộng hưởng là xác định — phát biểu đúng là: tốc độ truyền âm trong không khí. (Vì): Từ tần số đã biết và bước sóng đo qua cộng hưởng, ta tính tốc độ truyền âm. b) Sai: phương án nêu không phù hợp. c) Đúng: Âm thoa hoặc loa tạo âm có tần số biết trước, cần cho phép đo. d) Sai: phương án nêu không phù hợp
\item (Sai) Mục đích chính của thí nghiệm dùng ống cộng hưởng là xác định — phát biểu cường độ dòng điện là đúng.
\item (Đúng) Dụng cụ phát âm có tần số xác định trong thí nghiệm thường là — phát biểu đúng là: âm thoa hoặc loa nối máy phát âm tần.
\item (Sai) Dụng cụ phát âm có tần số xác định trong thí nghiệm thường là — phát biểu lực kế là đúng.
\end{itemchoice}
}
\end{ex}

% ===== Câu 62 =====
% ID: L11C2B15-16-TL
% Mức: VD
\dangbt{Tính tốc độ truyền âm từ số liệu thí nghiệm}
\begin{ex}
% ID: L11C2B15-16-TL
% Mức: VD
Trong dạng “Tính tốc độ truyền âm từ số liệu thí nghiệm”, Nêu hai biện pháp giảm sai số của phép đo.
\choiceTF

\loigiai{
Đặt mắt vuông góc khi đọc thước; xác định nhiều cặp cộng hưởng, đo lặp lại và lấy trung bình; giữ tần số, nhiệt độ ổn định.
}
\end{ex}

% ===== Câu 63 =====
% ID: L11C2B15-16-TLN
% Mức: VD
\begin{ex}
% ID: L11C2B15-16-TLN
% Mức: VD
Trong dạng “Tính tốc độ truyền âm từ số liệu thí nghiệm”, Với $\bar d=34$ cm và $f=500$ Hz, tính tốc độ âm theo m/s.
\shortans{340}
\loigiai{
$\lambda=2d=0,68$ m; $v=500\cdot0,68=340$ m/s.
}
\end{ex}

% ===== Câu 64 =====
% ID: L11C2B15-16-TN
% Mức: VD
\dangbt{Nhận biết mục đích và dụng cụ thí nghiệm đo tốc độ truyền âm}
\begin{ex}
% ID: L11C2B15-16-TN
% Mức: VD
Trong dạng “Nhận biết mục đích và dụng cụ thí nghiệm đo tốc độ truyền âm”, Để giảm sai số khi xác định khoảng cách cộng hưởng nên
\choice
{\True đo nhiều cặp vị trí và lấy trung bình}
{chỉ đo một lần}
{thay đổi tần số liên tục trong mỗi lần đo}
{đọc thước từ góc xiên}
\loigiai{
Đáp án A. Đo lặp lại và lấy trung bình làm giảm ảnh hưởng sai số ngẫu nhiên.
}
\end{ex}

% ===== Câu 65 =====
% ID: L11C2B15-17-DS
% Mức: TH
\dangbt{Xác định các vị trí cộng hưởng liên tiếp}
\begin{ex}
% ID: L11C2B15-17-DS
% Mức: TH
Xét các phát biểu sau trong dạng “Xác định các vị trí cộng hưởng liên tiếp”.
\choiceTF
{\True Khoảng cách giữa hai vị trí cộng hưởng liên tiếp trong ống một đầu kín xấp xỉ — phát biểu đúng là: $\lambda/2$.}
{Khoảng cách giữa hai vị trí cộng hưởng liên tiếp trong ống một đầu kín xấp xỉ — phát biểu $\lambda$ là đúng.}
{\True Biết khoảng cách hai vị trí cộng hưởng liên tiếp là d, bước sóng bằng — phát biểu đúng là: $2d$.}
{Biết khoảng cách hai vị trí cộng hưởng liên tiếp là d, bước sóng bằng — phát biểu $d$ là đúng.}
\loigiai{
\begin{itemchoice}
\item (Đúng) Khoảng cách giữa hai vị trí cộng hưởng liên tiếp trong ống một đầu kín xấp xỉ — phát biểu đúng là: $\lambda/2$. (Vì): Hai trạng thái cộng hưởng liên tiếp chênh nhau nửa bước sóng. b) Sai: phương án nêu không phù hợp. c) Đúng: Vì $d=\lambda/2$ nên $\lambda=2d$. d) Sai: phương án nêu không phù hợp
\item (Sai) Khoảng cách giữa hai vị trí cộng hưởng liên tiếp trong ống một đầu kín xấp xỉ — phát biểu $\lambda$ là đúng.
\item (Đúng) Biết khoảng cách hai vị trí cộng hưởng liên tiếp là d, bước sóng bằng — phát biểu đúng là: $2d$.
\item (Sai) Biết khoảng cách hai vị trí cộng hưởng liên tiếp là d, bước sóng bằng — phát biểu $d$ là đúng.
\end{itemchoice}
}
\end{ex}

% ===== Câu 66 =====
% ID: L11C2B15-17-TL
% Mức: VD
\dangbt{Tính tốc độ truyền âm từ số liệu thí nghiệm}
\begin{ex}
% ID: L11C2B15-17-TL
% Mức: VD
Trong dạng “Tính tốc độ truyền âm từ số liệu thí nghiệm”, Phân biệt đại lượng đo trực tiếp và gián tiếp trong thí nghiệm.
\choiceTF

\loigiai{
Chiều dài cột khí và tần số hiển thị là số liệu trực tiếp; bước sóng và tốc độ âm được suy ra bằng công thức nên là đại lượng gián tiếp.
}
\end{ex}

% ===== Câu 67 =====
% ID: L11C2B15-17-TLN
% Mức: VD
\begin{ex}
% ID: L11C2B15-17-TLN
% Mức: VD
Trong dạng “Tính tốc độ truyền âm từ số liệu thí nghiệm”, Giá trị đo là $338\,\text{m}$ /s, giá trị tham chiếu $340\,\text{m}$ /s. Tính độ lệch tuyệt đối theo m/s.
\shortans{2}
\loigiai{
$\Delta v=|338-340|=2$ m/s.
}
\end{ex}

% ===== Câu 68 =====
% ID: L11C2B15-17-TN
% Mức: VD
\dangbt{Nhận biết mục đích và dụng cụ thí nghiệm đo tốc độ truyền âm}
\begin{ex}
% ID: L11C2B15-17-TN
% Mức: VD
Trong dạng “Nhận biết mục đích và dụng cụ thí nghiệm đo tốc độ truyền âm”, Khi đọc thước, mắt phải đặt
\choice
{\True vuông góc với vạch chia tại vị trí cần đọc}
{càng xa thước càng tốt}
{xiên 45 độ}
{ở bất kì vị trí nào}
\loigiai{
Đáp án A. Đặt mắt vuông góc giúp tránh sai số thị sai.
}
\end{ex}

% ===== Câu 69 =====
% ID: L11C2B15-18-DS
% Mức: VD
\dangbt{Xác định các vị trí cộng hưởng liên tiếp}
\begin{ex}
% ID: L11C2B15-18-DS
% Mức: VD
Xét các phát biểu sau trong dạng “Xác định các vị trí cộng hưởng liên tiếp”.
\choiceTF
{\True Công thức tính tốc độ truyền âm là — phát biểu đúng là: $v=f\lambda$.}
{Công thức tính tốc độ truyền âm là — phát biểu $v=f/\lambda$ là đúng.}
{\True Để giảm sai số khi xác định khoảng cách cộng hưởng nên — phát biểu đúng là: đo nhiều cặp vị trí và lấy trung bình.}
{Để giảm sai số khi xác định khoảng cách cộng hưởng nên — phát biểu thay đổi tần số liên tục trong mỗi lần đo là đúng.}
\loigiai{
\begin{itemchoice}
\item (Đúng) Công thức tính tốc độ truyền âm là — phát biểu đúng là: $v=f\lambda$. (Vì): Tốc độ truyền sóng thỏa $v=f\lambda$. b) Sai: phương án nêu không phù hợp. c) Đúng: Đo lặp lại và lấy trung bình làm giảm ảnh hưởng sai số ngẫu nhiên. d) Sai: phương án nêu không phù hợp
\item (Sai) Công thức tính tốc độ truyền âm là — phát biểu $v=f/\lambda$ là đúng.
\item (Đúng) Để giảm sai số khi xác định khoảng cách cộng hưởng nên — phát biểu đúng là: đo nhiều cặp vị trí và lấy trung bình.
\item (Sai) Để giảm sai số khi xác định khoảng cách cộng hưởng nên — phát biểu thay đổi tần số liên tục trong mỗi lần đo là đúng.
\end{itemchoice}
}
\end{ex}

% ===== Câu 70 =====
% ID: L11C2B15-18-TL
% Mức: VDC
\dangbt{Tính tốc độ truyền âm từ số liệu thí nghiệm}
\begin{ex}
% ID: L11C2B15-18-TL
% Mức: VDC
Trong dạng “Tính tốc độ truyền âm từ số liệu thí nghiệm”, Một nhóm đo được $342\,\text{m}$ /s trong khi giá trị tham chiếu là $340\,\text{m}$ /s. Nhận xét kết quả.
\choiceTF

\loigiai{
Độ lệch tuyệt đối là $2\,\text{m}$ /s, sai lệch tương đối khoảng $2/340\cdot100\%=0,59\%$; kết quả khá phù hợp trong điều kiện thí nghiệm phổ thông.
}
\end{ex}

% ===== Câu 71 =====
% ID: L11C2B15-18-TLN
% Mức: VDC
\begin{ex}
% ID: L11C2B15-18-TLN
% Mức: VDC
Trong dạng “Tính tốc độ truyền âm từ số liệu thí nghiệm”, Đo được $d=0,30$ m với âm $600\,\text{Hz}$. Tính tốc độ âm theo m/s.
\shortans{360}
\loigiai{
$\lambda=2d=0,60$ m; $v=f\lambda=600\cdot0,60=360$ m/s.
}
\end{ex}

% ===== Câu 72 =====
% ID: L11C2B15-18-TN
% Mức: VD
\dangbt{Nhận biết mục đích và dụng cụ thí nghiệm đo tốc độ truyền âm}
\begin{ex}
% ID: L11C2B15-18-TN
% Mức: VD
Trong dạng “Nhận biết mục đích và dụng cụ thí nghiệm đo tốc độ truyền âm”, Nếu $f=500$ Hz và $\lambda=0,68$ m thì tốc độ âm là
\choice
{\True $340\,\text{m}$ /s}
{$735\,\text{m}$ /s}
{$0,00136\,\text{m}$ /s}
{$500,68\,\text{m}$ /s}
\loigiai{
Đáp án A. $v=f\lambda=500\cdot0,68=340$ m/s.
}
\end{ex}

% ===== Câu 73 =====
% ID: L11C2B15-19-DS
% Mức: VD
\dangbt{Xác định các vị trí cộng hưởng liên tiếp}
\begin{ex}
% ID: L11C2B15-19-DS
% Mức: VD
Xét các phát biểu sau trong dạng “Xác định các vị trí cộng hưởng liên tiếp”.
\choiceTF
{\True Khi đọc thước, mắt phải đặt — phát biểu đúng là: vuông góc với vạch chia tại vị trí cần đọc.}
{Khi đọc thước, mắt phải đặt — phát biểu càng xa thước càng tốt là đúng.}
{\True Nếu $f=500$ Hz và $\lambda=0,68$ m thì tốc độ âm là — phát biểu đúng là: $340\,\text{m}$ /s.}
{Nếu $f=500$ Hz và $\lambda=0,68$ m thì tốc độ âm là — phát biểu $0,00136\,\text{m}$ /s là đúng.}
\loigiai{
\begin{itemchoice}
\item (Đúng) Khi đọc thước, mắt phải đặt — phát biểu đúng là: vuông góc với vạch chia tại vị trí cần đọc. (Vì): ặt mắt vuông góc giúp tránh sai số thị sai. b) Sai: phương án nêu không phù hợp. c) Đúng: $v=f\lambda=500\cdot0,68=340$ m/s. d) Sai: phương án nêu không phù hợp
\item (Sai) Khi đọc thước, mắt phải đặt — phát biểu càng xa thước càng tốt là đúng.
\item (Đúng) Nếu $f=500$ Hz và $\lambda=0,68$ m thì tốc độ âm là — phát biểu đúng là: $340\,\text{m}$ /s.
\item (Sai) Nếu $f=500$ Hz và $\lambda=0,68$ m thì tốc độ âm là — phát biểu $0,00136\,\text{m}$ /s là đúng.
\end{itemchoice}
}
\end{ex}

% ===== Câu 74 =====
% ID: L11C2B15-19-TL
% Mức: TH
\begin{ex}
% ID: L11C2B15-19-TL
% Mức: TH
Trong dạng “Xác định các vị trí cộng hưởng liên tiếp”, Nêu các bước chính đo tốc độ truyền âm bằng ống cộng hưởng.
\choiceTF

\loigiai{
Phát âm có tần số f đã biết; thay đổi chiều dài cột khí; ghi hai vị trí cộng hưởng liên tiếp; tính $\lambda=2|L_2-L_1|$ rồi tính $v=f\lambda$.
}
\end{ex}

% ===== Câu 75 =====
% ID: L11C2B15-19-TLN
% Mức: TH
\begin{ex}
% ID: L11C2B15-19-TLN
% Mức: TH
Trong dạng “Xác định các vị trí cộng hưởng liên tiếp”, Hai vị trí cộng hưởng liên tiếp cách nhau $34\,\text{cm}$. Tính bước sóng theo cm.
\shortans{68}
\loigiai{
$\lambda=2d=2\cdot34=68$ cm.
}
\end{ex}

% ===== Câu 76 =====
% ID: L11C2B15-19-TN
% Mức: VDC
\dangbt{Nhận biết mục đích và dụng cụ thí nghiệm đo tốc độ truyền âm}
\begin{ex}
% ID: L11C2B15-19-TN
% Mức: VDC
Trong dạng “Nhận biết mục đích và dụng cụ thí nghiệm đo tốc độ truyền âm”, Trong thí nghiệm, đại lượng thường được đo trực tiếp là
\choice
{\True chiều dài cột khí tại cộng hưởng}
{tốc độ âm}
{bước sóng chính xác tuyệt đối}
{pha ban đầu của nguồn}
\loigiai{
Đáp án A. Ta đọc trực tiếp chiều dài cột khí; bước sóng và tốc độ được tính gián tiếp.
}
\end{ex}

% ===== Câu 77 =====
% ID: L11C2B15-20-DS
% Mức: VDC
\dangbt{Xác định các vị trí cộng hưởng liên tiếp}
\begin{ex}
% ID: L11C2B15-20-DS
% Mức: VDC
Xét các phát biểu sau trong dạng “Xác định các vị trí cộng hưởng liên tiếp”.
\choiceTF
{\True Trong thí nghiệm, đại lượng thường được đo trực tiếp là — phát biểu đúng là: chiều dài cột khí tại cộng hưởng.}
{Trong thí nghiệm, đại lượng thường được đo trực tiếp là — phát biểu tốc độ âm là đúng.}
{\True Khi nhiệt độ không khí tăng, tốc độ truyền âm trong không khí nhìn chung — phát biểu đúng là: tăng.}
{Khi nhiệt độ không khí tăng, tốc độ truyền âm trong không khí nhìn chung — phát biểu không đổi là đúng.}
\loigiai{
\begin{itemchoice}
\item (Đúng) Trong thí nghiệm, đại lượng thường được đo trực tiếp là — phát biểu đúng là: chiều dài cột khí tại cộng hưởng. (Vì): Ta đọc trực tiếp chiều dài cột khí; bước sóng và tốc độ được tính gián tiếp. b) Sai: phương án nêu không phù hợp. c) Đúng: Tốc độ âm trong khí tăng khi nhiệt độ tăng. d) Sai: phương án nêu không phù hợp
\item (Sai) Trong thí nghiệm, đại lượng thường được đo trực tiếp là — phát biểu tốc độ âm là đúng.
\item (Đúng) Khi nhiệt độ không khí tăng, tốc độ truyền âm trong không khí nhìn chung — phát biểu đúng là: tăng.
\item (Sai) Khi nhiệt độ không khí tăng, tốc độ truyền âm trong không khí nhìn chung — phát biểu không đổi là đúng.
\end{itemchoice}
}
\end{ex}

% ===== Câu 78 =====
% ID: L11C2B15-20-TL
% Mức: TH
\begin{ex}
% ID: L11C2B15-20-TL
% Mức: TH
Trong dạng “Xác định các vị trí cộng hưởng liên tiếp”, Giải thích vì sao dùng hai vị trí cộng hưởng liên tiếp giúp xác định bước sóng.
\choiceTF

\loigiai{
Trong ống một đầu kín, các chiều dài cộng hưởng kế tiếp chênh nhau $\lambda/2$, nên đo hiệu hai vị trí loại được phần hiệu chỉnh đầu ống và cho $\lambda=2\Delta L$.
}
\end{ex}

% ===== Câu 79 =====
% ID: L11C2B15-20-TLN
% Mức: TH
\begin{ex}
% ID: L11C2B15-20-TLN
% Mức: TH
Trong dạng “Xác định các vị trí cộng hưởng liên tiếp”, Âm có tần số $500\,\text{Hz}$, bước sóng $0,68\,\text{m}$. Tính tốc độ theo m/s.
\shortans{340}
\loigiai{
$v=f\lambda=500\cdot0,68=340$ m/s.
}
\end{ex}

% ===== Câu 80 =====
% ID: L11C2B15-20-TN
% Mức: VDC
\dangbt{Nhận biết mục đích và dụng cụ thí nghiệm đo tốc độ truyền âm}
\begin{ex}
% ID: L11C2B15-20-TN
% Mức: VDC
Trong dạng “Nhận biết mục đích và dụng cụ thí nghiệm đo tốc độ truyền âm”, Khi nhiệt độ không khí tăng, tốc độ truyền âm trong không khí nhìn chung
\choice
{\True tăng}
{giảm}
{không đổi}
{bằng 0}
\loigiai{
Đáp án A. Tốc độ âm trong khí tăng khi nhiệt độ tăng.
}
\end{ex}

% ===== Câu 81 =====
% ID: L11C2B15-21-DS
% Mức: TH
\dangbt{Xử lí số liệu và đánh giá sai số phép đo}
\begin{ex}
% ID: L11C2B15-21-DS
% Mức: TH
Xét các phát biểu sau trong dạng “Xử lí số liệu và đánh giá sai số phép đo”.
\choiceTF
{\True Mục đích chính của thí nghiệm dùng ống cộng hưởng là xác định — phát biểu đúng là: tốc độ truyền âm trong không khí.}
{Mục đích chính của thí nghiệm dùng ống cộng hưởng là xác định — phát biểu cường độ dòng điện là đúng.}
{\True Dụng cụ phát âm có tần số xác định trong thí nghiệm thường là — phát biểu đúng là: âm thoa hoặc loa nối máy phát âm tần.}
{Dụng cụ phát âm có tần số xác định trong thí nghiệm thường là — phát biểu lực kế là đúng.}
\loigiai{
\begin{itemchoice}
\item (Đúng) Mục đích chính của thí nghiệm dùng ống cộng hưởng là xác định — phát biểu đúng là: tốc độ truyền âm trong không khí. (Vì): Từ tần số đã biết và bước sóng đo qua cộng hưởng, ta tính tốc độ truyền âm. b) Sai: phương án nêu không phù hợp. c) Đúng: Âm thoa hoặc loa tạo âm có tần số biết trước, cần cho phép đo. d) Sai: phương án nêu không phù hợp
\item (Sai) Mục đích chính của thí nghiệm dùng ống cộng hưởng là xác định — phát biểu cường độ dòng điện là đúng.
\item (Đúng) Dụng cụ phát âm có tần số xác định trong thí nghiệm thường là — phát biểu đúng là: âm thoa hoặc loa nối máy phát âm tần.
\item (Sai) Dụng cụ phát âm có tần số xác định trong thí nghiệm thường là — phát biểu lực kế là đúng.
\end{itemchoice}
}
\end{ex}

% ===== Câu 82 =====
% ID: L11C2B15-21-TL
% Mức: VD
\dangbt{Xác định các vị trí cộng hưởng liên tiếp}
\begin{ex}
% ID: L11C2B15-21-TL
% Mức: VD
Trong dạng “Xác định các vị trí cộng hưởng liên tiếp”, Với $f=400$ Hz, hai vị trí cộng hưởng là $18\,\text{cm}$ và $60\,\text{cm}$. Tính tốc độ âm.
\choiceTF

\loigiai{
$d=60-18=42\,\mathrm{cm}=0,42$ m; $\lambda=0,84$ m; $v=400\cdot0,84=336$ m/s.
}
\end{ex}

% ===== Câu 83 =====
% ID: L11C2B15-21-TLN
% Mức: VD
\begin{ex}
% ID: L11C2B15-21-TLN
% Mức: VD
Trong dạng “Xác định các vị trí cộng hưởng liên tiếp”, Các khoảng cách cộng hưởng đo được là $33,8\,\text{cm}$; $34,0\,\text{cm}$; $34,2\,\text{cm}$. Tính giá trị trung bình theo cm.
\shortans{34}
\loigiai{
$\bar d=(33,8+34,0+34,2)/3=34,0$ cm.
}
\end{ex}

% ===== Câu 84 =====
% ID: L11C2B15-21-TN
% Mức: NB
\dangbt{Tính tốc độ truyền âm từ số liệu thí nghiệm}
\begin{ex}
% ID: L11C2B15-21-TN
% Mức: NB
Trong dạng “Tính tốc độ truyền âm từ số liệu thí nghiệm”, Mục đích chính của thí nghiệm dùng ống cộng hưởng là xác định
\choice
{\True tốc độ truyền âm trong không khí}
{cường độ dòng điện}
{khối lượng riêng của nước}
{nhiệt dung riêng của không khí}
\loigiai{
Đáp án A. Từ tần số đã biết và bước sóng đo qua cộng hưởng, ta tính tốc độ truyền âm.
}
\end{ex}

% ===== Câu 85 =====
% ID: L11C2B15-22-DS
% Mức: TH
\dangbt{Xử lí số liệu và đánh giá sai số phép đo}
\begin{ex}
% ID: L11C2B15-22-DS
% Mức: TH
Xét các phát biểu sau trong dạng “Xử lí số liệu và đánh giá sai số phép đo”.
\choiceTF
{\True Khoảng cách giữa hai vị trí cộng hưởng liên tiếp trong ống một đầu kín xấp xỉ — phát biểu đúng là: $\lambda/2$.}
{Khoảng cách giữa hai vị trí cộng hưởng liên tiếp trong ống một đầu kín xấp xỉ — phát biểu $\lambda$ là đúng.}
{\True Biết khoảng cách hai vị trí cộng hưởng liên tiếp là d, bước sóng bằng — phát biểu đúng là: $2d$.}
{Biết khoảng cách hai vị trí cộng hưởng liên tiếp là d, bước sóng bằng — phát biểu $d$ là đúng.}
\loigiai{
\begin{itemchoice}
\item (Đúng) Khoảng cách giữa hai vị trí cộng hưởng liên tiếp trong ống một đầu kín xấp xỉ — phát biểu đúng là: $\lambda/2$. (Vì): Hai trạng thái cộng hưởng liên tiếp chênh nhau nửa bước sóng. b) Sai: phương án nêu không phù hợp. c) Đúng: Vì $d=\lambda/2$ nên $\lambda=2d$. d) Sai: phương án nêu không phù hợp
\item (Sai) Khoảng cách giữa hai vị trí cộng hưởng liên tiếp trong ống một đầu kín xấp xỉ — phát biểu $\lambda$ là đúng.
\item (Đúng) Biết khoảng cách hai vị trí cộng hưởng liên tiếp là d, bước sóng bằng — phát biểu đúng là: $2d$.
\item (Sai) Biết khoảng cách hai vị trí cộng hưởng liên tiếp là d, bước sóng bằng — phát biểu $d$ là đúng.
\end{itemchoice}
}
\end{ex}

% ===== Câu 86 =====
% ID: L11C2B15-22-TL
% Mức: VD
\dangbt{Xác định các vị trí cộng hưởng liên tiếp}
\begin{ex}
% ID: L11C2B15-22-TL
% Mức: VD
Trong dạng “Xác định các vị trí cộng hưởng liên tiếp”, Nêu hai biện pháp giảm sai số của phép đo.
\choiceTF

\loigiai{
Đặt mắt vuông góc khi đọc thước; xác định nhiều cặp cộng hưởng, đo lặp lại và lấy trung bình; giữ tần số, nhiệt độ ổn định.
}
\end{ex}

% ===== Câu 87 =====
% ID: L11C2B15-22-TLN
% Mức: VD
\begin{ex}
% ID: L11C2B15-22-TLN
% Mức: VD
Trong dạng “Xác định các vị trí cộng hưởng liên tiếp”, Với $\bar d=34$ cm và $f=500$ Hz, tính tốc độ âm theo m/s.
\shortans{340}
\loigiai{
$\lambda=2d=0,68$ m; $v=500\cdot0,68=340$ m/s.
}
\end{ex}

% ===== Câu 88 =====
% ID: L11C2B15-22-TN
% Mức: NB
\dangbt{Tính tốc độ truyền âm từ số liệu thí nghiệm}
\begin{ex}
% ID: L11C2B15-22-TN
% Mức: NB
Trong dạng “Tính tốc độ truyền âm từ số liệu thí nghiệm”, Dụng cụ phát âm có tần số xác định trong thí nghiệm thường là
\choice
{\True âm thoa hoặc loa nối máy phát âm tần}
{nhiệt kế}
{lực kế}
{cân điện tử}
\loigiai{
Đáp án A. Âm thoa hoặc loa tạo âm có tần số biết trước, cần cho phép đo.
}
\end{ex}

% ===== Câu 89 =====
% ID: L11C2B15-23-DS
% Mức: VD
\dangbt{Xử lí số liệu và đánh giá sai số phép đo}
\begin{ex}
% ID: L11C2B15-23-DS
% Mức: VD
Xét các phát biểu sau trong dạng “Xử lí số liệu và đánh giá sai số phép đo”.
\choiceTF
{\True Công thức tính tốc độ truyền âm là — phát biểu đúng là: $v=f\lambda$.}
{Công thức tính tốc độ truyền âm là — phát biểu $v=f/\lambda$ là đúng.}
{\True Để giảm sai số khi xác định khoảng cách cộng hưởng nên — phát biểu đúng là: đo nhiều cặp vị trí và lấy trung bình.}
{Để giảm sai số khi xác định khoảng cách cộng hưởng nên — phát biểu thay đổi tần số liên tục trong mỗi lần đo là đúng.}
\loigiai{
\begin{itemchoice}
\item (Đúng) Công thức tính tốc độ truyền âm là — phát biểu đúng là: $v=f\lambda$. (Vì): Tốc độ truyền sóng thỏa $v=f\lambda$. b) Sai: phương án nêu không phù hợp. c) Đúng: Đo lặp lại và lấy trung bình làm giảm ảnh hưởng sai số ngẫu nhiên. d) Sai: phương án nêu không phù hợp
\item (Sai) Công thức tính tốc độ truyền âm là — phát biểu $v=f/\lambda$ là đúng.
\item (Đúng) Để giảm sai số khi xác định khoảng cách cộng hưởng nên — phát biểu đúng là: đo nhiều cặp vị trí và lấy trung bình.
\item (Sai) Để giảm sai số khi xác định khoảng cách cộng hưởng nên — phát biểu thay đổi tần số liên tục trong mỗi lần đo là đúng.
\end{itemchoice}
}
\end{ex}

% ===== Câu 90 =====
% ID: L11C2B15-23-TL
% Mức: VD
\dangbt{Xác định các vị trí cộng hưởng liên tiếp}
\begin{ex}
% ID: L11C2B15-23-TL
% Mức: VD
Trong dạng “Xác định các vị trí cộng hưởng liên tiếp”, Phân biệt đại lượng đo trực tiếp và gián tiếp trong thí nghiệm.
\choiceTF

\loigiai{
Chiều dài cột khí và tần số hiển thị là số liệu trực tiếp; bước sóng và tốc độ âm được suy ra bằng công thức nên là đại lượng gián tiếp.
}
\end{ex}

% ===== Câu 91 =====
% ID: L11C2B15-23-TLN
% Mức: VD
\begin{ex}
% ID: L11C2B15-23-TLN
% Mức: VD
Trong dạng “Xác định các vị trí cộng hưởng liên tiếp”, Giá trị đo là $338\,\text{m}$ /s, giá trị tham chiếu $340\,\text{m}$ /s. Tính độ lệch tuyệt đối theo m/s.
\shortans{2}
\loigiai{
$\Delta v=|338-340|=2$ m/s.
}
\end{ex}

% ===== Câu 92 =====
% ID: L11C2B15-23-TN
% Mức: TH
\dangbt{Tính tốc độ truyền âm từ số liệu thí nghiệm}
\begin{ex}
% ID: L11C2B15-23-TN
% Mức: TH
Trong dạng “Tính tốc độ truyền âm từ số liệu thí nghiệm”, Khoảng cách giữa hai vị trí cộng hưởng liên tiếp trong ống một đầu kín xấp xỉ
\choice
{\True $\lambda/2$}
{$\lambda$}
{$\lambda/4$}
{$2\lambda$}
\loigiai{
Đáp án A. Hai trạng thái cộng hưởng liên tiếp chênh nhau nửa bước sóng.
}
\end{ex}

% ===== Câu 93 =====
% ID: L11C2B15-24-DS
% Mức: VD
\dangbt{Xử lí số liệu và đánh giá sai số phép đo}
\begin{ex}
% ID: L11C2B15-24-DS
% Mức: VD
Xét các phát biểu sau trong dạng “Xử lí số liệu và đánh giá sai số phép đo”.
\choiceTF
{\True Khi đọc thước, mắt phải đặt — phát biểu đúng là: vuông góc với vạch chia tại vị trí cần đọc.}
{Khi đọc thước, mắt phải đặt — phát biểu càng xa thước càng tốt là đúng.}
{\True Nếu $f=500$ Hz và $\lambda=0,68$ m thì tốc độ âm là — phát biểu đúng là: $340\,\text{m}$ /s.}
{Nếu $f=500$ Hz và $\lambda=0,68$ m thì tốc độ âm là — phát biểu $0,00136\,\text{m}$ /s là đúng.}
\loigiai{
\begin{itemchoice}
\item (Đúng) Khi đọc thước, mắt phải đặt — phát biểu đúng là: vuông góc với vạch chia tại vị trí cần đọc. (Vì): ặt mắt vuông góc giúp tránh sai số thị sai. b) Sai: phương án nêu không phù hợp. c) Đúng: $v=f\lambda=500\cdot0,68=340$ m/s. d) Sai: phương án nêu không phù hợp
\item (Sai) Khi đọc thước, mắt phải đặt — phát biểu càng xa thước càng tốt là đúng.
\item (Đúng) Nếu $f=500$ Hz và $\lambda=0,68$ m thì tốc độ âm là — phát biểu đúng là: $340\,\text{m}$ /s.
\item (Sai) Nếu $f=500$ Hz và $\lambda=0,68$ m thì tốc độ âm là — phát biểu $0,00136\,\text{m}$ /s là đúng.
\end{itemchoice}
}
\end{ex}

% ===== Câu 94 =====
% ID: L11C2B15-24-TL
% Mức: VDC
\dangbt{Xác định các vị trí cộng hưởng liên tiếp}
\begin{ex}
% ID: L11C2B15-24-TL
% Mức: VDC
Trong dạng “Xác định các vị trí cộng hưởng liên tiếp”, Một nhóm đo được $342\,\text{m}$ /s trong khi giá trị tham chiếu là $340\,\text{m}$ /s. Nhận xét kết quả.
\choiceTF

\loigiai{
Độ lệch tuyệt đối là $2\,\text{m}$ /s, sai lệch tương đối khoảng $2/340\cdot100\%=0,59\%$; kết quả khá phù hợp trong điều kiện thí nghiệm phổ thông.
}
\end{ex}

% ===== Câu 95 =====
% ID: L11C2B15-24-TLN
% Mức: VDC
\begin{ex}
% ID: L11C2B15-24-TLN
% Mức: VDC
Trong dạng “Xác định các vị trí cộng hưởng liên tiếp”, Đo được $d=0,30$ m với âm $600\,\text{Hz}$. Tính tốc độ âm theo m/s.
\shortans{360}
\loigiai{
$\lambda=2d=0,60$ m; $v=f\lambda=600\cdot0,60=360$ m/s.
}
\end{ex}

% ===== Câu 96 =====
% ID: L11C2B15-24-TN
% Mức: TH
\dangbt{Tính tốc độ truyền âm từ số liệu thí nghiệm}
\begin{ex}
% ID: L11C2B15-24-TN
% Mức: TH
Trong dạng “Tính tốc độ truyền âm từ số liệu thí nghiệm”, Biết khoảng cách hai vị trí cộng hưởng liên tiếp là d, bước sóng bằng
\choice
{\True $2d$}
{$d/2$}
{$d$}
{$4d$}
\loigiai{
Đáp án A. Vì $d=\lambda/2$ nên $\lambda=2d$.
}
\end{ex}

% ===== Câu 97 =====
% ID: L11C2B15-25-DS
% Mức: VDC
\dangbt{Xử lí số liệu và đánh giá sai số phép đo}
\begin{ex}
% ID: L11C2B15-25-DS
% Mức: VDC
Xét các phát biểu sau trong dạng “Xử lí số liệu và đánh giá sai số phép đo”.
\choiceTF
{\True Trong thí nghiệm, đại lượng thường được đo trực tiếp là — phát biểu đúng là: chiều dài cột khí tại cộng hưởng.}
{Trong thí nghiệm, đại lượng thường được đo trực tiếp là — phát biểu tốc độ âm là đúng.}
{\True Khi nhiệt độ không khí tăng, tốc độ truyền âm trong không khí nhìn chung — phát biểu đúng là: tăng.}
{Khi nhiệt độ không khí tăng, tốc độ truyền âm trong không khí nhìn chung — phát biểu không đổi là đúng.}
\loigiai{
\begin{itemchoice}
\item (Đúng) Trong thí nghiệm, đại lượng thường được đo trực tiếp là — phát biểu đúng là: chiều dài cột khí tại cộng hưởng. (Vì): Ta đọc trực tiếp chiều dài cột khí; bước sóng và tốc độ được tính gián tiếp. b) Sai: phương án nêu không phù hợp. c) Đúng: Tốc độ âm trong khí tăng khi nhiệt độ tăng. d) Sai: phương án nêu không phù hợp
\item (Sai) Trong thí nghiệm, đại lượng thường được đo trực tiếp là — phát biểu tốc độ âm là đúng.
\item (Đúng) Khi nhiệt độ không khí tăng, tốc độ truyền âm trong không khí nhìn chung — phát biểu đúng là: tăng.
\item (Sai) Khi nhiệt độ không khí tăng, tốc độ truyền âm trong không khí nhìn chung — phát biểu không đổi là đúng.
\end{itemchoice}
}
\end{ex}

% ===== Câu 98 =====
% ID: L11C2B15-25-TL
% Mức: TH
\begin{ex}
% ID: L11C2B15-25-TL
% Mức: TH
Trong dạng “Xử lí số liệu và đánh giá sai số phép đo”, Nêu các bước chính đo tốc độ truyền âm bằng ống cộng hưởng.
\choiceTF

\loigiai{
Phát âm có tần số f đã biết; thay đổi chiều dài cột khí; ghi hai vị trí cộng hưởng liên tiếp; tính $\lambda=2|L_2-L_1|$ rồi tính $v=f\lambda$.
}
\end{ex}

% ===== Câu 99 =====
% ID: L11C2B15-25-TLN
% Mức: TH
\begin{ex}
% ID: L11C2B15-25-TLN
% Mức: TH
Trong dạng “Xử lí số liệu và đánh giá sai số phép đo”, Hai vị trí cộng hưởng liên tiếp cách nhau $34\,\text{cm}$. Tính bước sóng theo cm.
\shortans{68}
\loigiai{
$\lambda=2d=2\cdot34=68$ cm.
}
\end{ex}

% ===== Câu 100 =====
% ID: L11C2B15-25-TN
% Mức: TH
\dangbt{Tính tốc độ truyền âm từ số liệu thí nghiệm}
\begin{ex}
% ID: L11C2B15-25-TN
% Mức: TH
Trong dạng “Tính tốc độ truyền âm từ số liệu thí nghiệm”, Công thức tính tốc độ truyền âm là
\choice
{\True $v=f\lambda$}
{$v=f/\lambda$}
{$v=\lambda/f$}
{$v=f+\lambda$}
\loigiai{
Đáp án A. Tốc độ truyền sóng thỏa $v=f\lambda$.
}
\end{ex}

% ===== Câu 101 =====
% ID: L11C2B15-26-DS
% Mức: TH
\dangbt{Đề xuất cải tiến và vận dụng thí nghiệm đo tốc độ truyền âm}
\begin{ex}
% ID: L11C2B15-26-DS
% Mức: TH
Xét các phát biểu sau trong dạng “Đề xuất cải tiến và vận dụng thí nghiệm đo tốc độ truyền âm”.
\choiceTF
{\True Mục đích chính của thí nghiệm dùng ống cộng hưởng là xác định — phát biểu đúng là: tốc độ truyền âm trong không khí.}
{Mục đích chính của thí nghiệm dùng ống cộng hưởng là xác định — phát biểu cường độ dòng điện là đúng.}
{\True Dụng cụ phát âm có tần số xác định trong thí nghiệm thường là — phát biểu đúng là: âm thoa hoặc loa nối máy phát âm tần.}
{Dụng cụ phát âm có tần số xác định trong thí nghiệm thường là — phát biểu lực kế là đúng.}
\loigiai{
\begin{itemchoice}
\item (Đúng) Mục đích chính của thí nghiệm dùng ống cộng hưởng là xác định — phát biểu đúng là: tốc độ truyền âm trong không khí. (Vì): Từ tần số đã biết và bước sóng đo qua cộng hưởng, ta tính tốc độ truyền âm. b) Sai: phương án nêu không phù hợp. c) Đúng: Âm thoa hoặc loa tạo âm có tần số biết trước, cần cho phép đo. d) Sai: phương án nêu không phù hợp
\item (Sai) Mục đích chính của thí nghiệm dùng ống cộng hưởng là xác định — phát biểu cường độ dòng điện là đúng.
\item (Đúng) Dụng cụ phát âm có tần số xác định trong thí nghiệm thường là — phát biểu đúng là: âm thoa hoặc loa nối máy phát âm tần.
\item (Sai) Dụng cụ phát âm có tần số xác định trong thí nghiệm thường là — phát biểu lực kế là đúng.
\end{itemchoice}
}
\end{ex}

% ===== Câu 102 =====
% ID: L11C2B15-26-TL
% Mức: TH
\dangbt{Xử lí số liệu và đánh giá sai số phép đo}
\begin{ex}
% ID: L11C2B15-26-TL
% Mức: TH
Trong dạng “Xử lí số liệu và đánh giá sai số phép đo”, Giải thích vì sao dùng hai vị trí cộng hưởng liên tiếp giúp xác định bước sóng.
\choiceTF

\loigiai{
Trong ống một đầu kín, các chiều dài cộng hưởng kế tiếp chênh nhau $\lambda/2$, nên đo hiệu hai vị trí loại được phần hiệu chỉnh đầu ống và cho $\lambda=2\Delta L$.
}
\end{ex}

% ===== Câu 103 =====
% ID: L11C2B15-26-TLN
% Mức: TH
\begin{ex}
% ID: L11C2B15-26-TLN
% Mức: TH
Trong dạng “Xử lí số liệu và đánh giá sai số phép đo”, Âm có tần số $500\,\text{Hz}$, bước sóng $0,68\,\text{m}$. Tính tốc độ theo m/s.
\shortans{340}
\loigiai{
$v=f\lambda=500\cdot0,68=340$ m/s.
}
\end{ex}

% ===== Câu 104 =====
% ID: L11C2B15-26-TN
% Mức: VD
\dangbt{Tính tốc độ truyền âm từ số liệu thí nghiệm}
\begin{ex}
% ID: L11C2B15-26-TN
% Mức: VD
Trong dạng “Tính tốc độ truyền âm từ số liệu thí nghiệm”, Để giảm sai số khi xác định khoảng cách cộng hưởng nên
\choice
{\True đo nhiều cặp vị trí và lấy trung bình}
{chỉ đo một lần}
{thay đổi tần số liên tục trong mỗi lần đo}
{đọc thước từ góc xiên}
\loigiai{
Đáp án A. Đo lặp lại và lấy trung bình làm giảm ảnh hưởng sai số ngẫu nhiên.
}
\end{ex}

% ===== Câu 105 =====
% ID: L11C2B15-27-DS
% Mức: TH
\dangbt{Đề xuất cải tiến và vận dụng thí nghiệm đo tốc độ truyền âm}
\begin{ex}
% ID: L11C2B15-27-DS
% Mức: TH
Xét các phát biểu sau trong dạng “Đề xuất cải tiến và vận dụng thí nghiệm đo tốc độ truyền âm”.
\choiceTF
{\True Khoảng cách giữa hai vị trí cộng hưởng liên tiếp trong ống một đầu kín xấp xỉ — phát biểu đúng là: $\lambda/2$.}
{Khoảng cách giữa hai vị trí cộng hưởng liên tiếp trong ống một đầu kín xấp xỉ — phát biểu $\lambda$ là đúng.}
{\True Biết khoảng cách hai vị trí cộng hưởng liên tiếp là d, bước sóng bằng — phát biểu đúng là: $2d$.}
{Biết khoảng cách hai vị trí cộng hưởng liên tiếp là d, bước sóng bằng — phát biểu $d$ là đúng.}
\loigiai{
\begin{itemchoice}
\item (Đúng) Khoảng cách giữa hai vị trí cộng hưởng liên tiếp trong ống một đầu kín xấp xỉ — phát biểu đúng là: $\lambda/2$. (Vì): Hai trạng thái cộng hưởng liên tiếp chênh nhau nửa bước sóng. b) Sai: phương án nêu không phù hợp. c) Đúng: Vì $d=\lambda/2$ nên $\lambda=2d$. d) Sai: phương án nêu không phù hợp
\item (Sai) Khoảng cách giữa hai vị trí cộng hưởng liên tiếp trong ống một đầu kín xấp xỉ — phát biểu $\lambda$ là đúng.
\item (Đúng) Biết khoảng cách hai vị trí cộng hưởng liên tiếp là d, bước sóng bằng — phát biểu đúng là: $2d$.
\item (Sai) Biết khoảng cách hai vị trí cộng hưởng liên tiếp là d, bước sóng bằng — phát biểu $d$ là đúng.
\end{itemchoice}
}
\end{ex}

% ===== Câu 106 =====
% ID: L11C2B15-27-TL
% Mức: VD
\dangbt{Xử lí số liệu và đánh giá sai số phép đo}
\begin{ex}
% ID: L11C2B15-27-TL
% Mức: VD
Trong dạng “Xử lí số liệu và đánh giá sai số phép đo”, Với $f=400$ Hz, hai vị trí cộng hưởng là $18\,\text{cm}$ và $60\,\text{cm}$. Tính tốc độ âm.
\choiceTF

\loigiai{
$d=60-18=42\,\mathrm{cm}=0,42$ m; $\lambda=0,84$ m; $v=400\cdot0,84=336$ m/s.
}
\end{ex}

% ===== Câu 107 =====
% ID: L11C2B15-27-TLN
% Mức: VD
\begin{ex}
% ID: L11C2B15-27-TLN
% Mức: VD
Trong dạng “Xử lí số liệu và đánh giá sai số phép đo”, Các khoảng cách cộng hưởng đo được là $33,8\,\text{cm}$; $34,0\,\text{cm}$; $34,2\,\text{cm}$. Tính giá trị trung bình theo cm.
\shortans{34}
\loigiai{
$\bar d=(33,8+34,0+34,2)/3=34,0$ cm.
}
\end{ex}

% ===== Câu 108 =====
% ID: L11C2B15-27-TN
% Mức: VD
\dangbt{Tính tốc độ truyền âm từ số liệu thí nghiệm}
\begin{ex}
% ID: L11C2B15-27-TN
% Mức: VD
Trong dạng “Tính tốc độ truyền âm từ số liệu thí nghiệm”, Khi đọc thước, mắt phải đặt
\choice
{\True vuông góc với vạch chia tại vị trí cần đọc}
{càng xa thước càng tốt}
{xiên 45 độ}
{ở bất kì vị trí nào}
\loigiai{
Đáp án A. Đặt mắt vuông góc giúp tránh sai số thị sai.
}
\end{ex}

% ===== Câu 109 =====
% ID: L11C2B15-28-DS
% Mức: VD
\dangbt{Đề xuất cải tiến và vận dụng thí nghiệm đo tốc độ truyền âm}
\begin{ex}
% ID: L11C2B15-28-DS
% Mức: VD
Xét các phát biểu sau trong dạng “Đề xuất cải tiến và vận dụng thí nghiệm đo tốc độ truyền âm”.
\choiceTF
{\True Công thức tính tốc độ truyền âm là — phát biểu đúng là: $v=f\lambda$.}
{Công thức tính tốc độ truyền âm là — phát biểu $v=f/\lambda$ là đúng.}
{\True Để giảm sai số khi xác định khoảng cách cộng hưởng nên — phát biểu đúng là: đo nhiều cặp vị trí và lấy trung bình.}
{Để giảm sai số khi xác định khoảng cách cộng hưởng nên — phát biểu thay đổi tần số liên tục trong mỗi lần đo là đúng.}
\loigiai{
\begin{itemchoice}
\item (Đúng) Công thức tính tốc độ truyền âm là — phát biểu đúng là: $v=f\lambda$. (Vì): Tốc độ truyền sóng thỏa $v=f\lambda$. b) Sai: phương án nêu không phù hợp. c) Đúng: Đo lặp lại và lấy trung bình làm giảm ảnh hưởng sai số ngẫu nhiên. d) Sai: phương án nêu không phù hợp
\item (Sai) Công thức tính tốc độ truyền âm là — phát biểu $v=f/\lambda$ là đúng.
\item (Đúng) Để giảm sai số khi xác định khoảng cách cộng hưởng nên — phát biểu đúng là: đo nhiều cặp vị trí và lấy trung bình.
\item (Sai) Để giảm sai số khi xác định khoảng cách cộng hưởng nên — phát biểu thay đổi tần số liên tục trong mỗi lần đo là đúng.
\end{itemchoice}
}
\end{ex}

% ===== Câu 110 =====
% ID: L11C2B15-28-TL
% Mức: VD
\dangbt{Xử lí số liệu và đánh giá sai số phép đo}
\begin{ex}
% ID: L11C2B15-28-TL
% Mức: VD
Trong dạng “Xử lí số liệu và đánh giá sai số phép đo”, Nêu hai biện pháp giảm sai số của phép đo.
\choiceTF

\loigiai{
Đặt mắt vuông góc khi đọc thước; xác định nhiều cặp cộng hưởng, đo lặp lại và lấy trung bình; giữ tần số, nhiệt độ ổn định.
}
\end{ex}

% ===== Câu 111 =====
% ID: L11C2B15-28-TLN
% Mức: VD
\begin{ex}
% ID: L11C2B15-28-TLN
% Mức: VD
Trong dạng “Xử lí số liệu và đánh giá sai số phép đo”, Với $\bar d=34$ cm và $f=500$ Hz, tính tốc độ âm theo m/s.
\shortans{340}
\loigiai{
$\lambda=2d=0,68$ m; $v=500\cdot0,68=340$ m/s.
}
\end{ex}

% ===== Câu 112 =====
% ID: L11C2B15-28-TN
% Mức: VD
\dangbt{Tính tốc độ truyền âm từ số liệu thí nghiệm}
\begin{ex}
% ID: L11C2B15-28-TN
% Mức: VD
Trong dạng “Tính tốc độ truyền âm từ số liệu thí nghiệm”, Nếu $f=500$ Hz và $\lambda=0,68$ m thì tốc độ âm là
\choice
{\True $340\,\text{m}$ /s}
{$735\,\text{m}$ /s}
{$0,00136\,\text{m}$ /s}
{$500,68\,\text{m}$ /s}
\loigiai{
Đáp án A. $v=f\lambda=500\cdot0,68=340$ m/s.
}
\end{ex}

% ===== Câu 113 =====
% ID: L11C2B15-29-DS
% Mức: VD
\dangbt{Đề xuất cải tiến và vận dụng thí nghiệm đo tốc độ truyền âm}
\begin{ex}
% ID: L11C2B15-29-DS
% Mức: VD
Xét các phát biểu sau trong dạng “Đề xuất cải tiến và vận dụng thí nghiệm đo tốc độ truyền âm”.
\choiceTF
{\True Khi đọc thước, mắt phải đặt — phát biểu đúng là: vuông góc với vạch chia tại vị trí cần đọc.}
{Khi đọc thước, mắt phải đặt — phát biểu càng xa thước càng tốt là đúng.}
{\True Nếu $f=500$ Hz và $\lambda=0,68$ m thì tốc độ âm là — phát biểu đúng là: $340\,\text{m}$ /s.}
{Nếu $f=500$ Hz và $\lambda=0,68$ m thì tốc độ âm là — phát biểu $0,00136\,\text{m}$ /s là đúng.}
\loigiai{
\begin{itemchoice}
\item (Đúng) Khi đọc thước, mắt phải đặt — phát biểu đúng là: vuông góc với vạch chia tại vị trí cần đọc. (Vì): ặt mắt vuông góc giúp tránh sai số thị sai. b) Sai: phương án nêu không phù hợp. c) Đúng: $v=f\lambda=500\cdot0,68=340$ m/s. d) Sai: phương án nêu không phù hợp
\item (Sai) Khi đọc thước, mắt phải đặt — phát biểu càng xa thước càng tốt là đúng.
\item (Đúng) Nếu $f=500$ Hz và $\lambda=0,68$ m thì tốc độ âm là — phát biểu đúng là: $340\,\text{m}$ /s.
\item (Sai) Nếu $f=500$ Hz và $\lambda=0,68$ m thì tốc độ âm là — phát biểu $0,00136\,\text{m}$ /s là đúng.
\end{itemchoice}
}
\end{ex}

% ===== Câu 114 =====
% ID: L11C2B15-29-TL
% Mức: VD
\dangbt{Xử lí số liệu và đánh giá sai số phép đo}
\begin{ex}
% ID: L11C2B15-29-TL
% Mức: VD
Trong dạng “Xử lí số liệu và đánh giá sai số phép đo”, Phân biệt đại lượng đo trực tiếp và gián tiếp trong thí nghiệm.
\choiceTF

\loigiai{
Chiều dài cột khí và tần số hiển thị là số liệu trực tiếp; bước sóng và tốc độ âm được suy ra bằng công thức nên là đại lượng gián tiếp.
}
\end{ex}

% ===== Câu 115 =====
% ID: L11C2B15-29-TLN
% Mức: VD
\begin{ex}
% ID: L11C2B15-29-TLN
% Mức: VD
Trong dạng “Xử lí số liệu và đánh giá sai số phép đo”, Giá trị đo là $338\,\text{m}$ /s, giá trị tham chiếu $340\,\text{m}$ /s. Tính độ lệch tuyệt đối theo m/s.
\shortans{2}
\loigiai{
$\Delta v=|338-340|=2$ m/s.
}
\end{ex}

% ===== Câu 116 =====
% ID: L11C2B15-29-TN
% Mức: VDC
\dangbt{Tính tốc độ truyền âm từ số liệu thí nghiệm}
\begin{ex}
% ID: L11C2B15-29-TN
% Mức: VDC
Trong dạng “Tính tốc độ truyền âm từ số liệu thí nghiệm”, Trong thí nghiệm, đại lượng thường được đo trực tiếp là
\choice
{\True chiều dài cột khí tại cộng hưởng}
{tốc độ âm}
{bước sóng chính xác tuyệt đối}
{pha ban đầu của nguồn}
\loigiai{
Đáp án A. Ta đọc trực tiếp chiều dài cột khí; bước sóng và tốc độ được tính gián tiếp.
}
\end{ex}

% ===== Câu 117 =====
% ID: L11C2B15-30-DS
% Mức: VDC
\dangbt{Đề xuất cải tiến và vận dụng thí nghiệm đo tốc độ truyền âm}
\begin{ex}
% ID: L11C2B15-30-DS
% Mức: VDC
Xét các phát biểu sau trong dạng “Đề xuất cải tiến và vận dụng thí nghiệm đo tốc độ truyền âm”.
\choiceTF
{\True Trong thí nghiệm, đại lượng thường được đo trực tiếp là — phát biểu đúng là: chiều dài cột khí tại cộng hưởng.}
{Trong thí nghiệm, đại lượng thường được đo trực tiếp là — phát biểu tốc độ âm là đúng.}
{\True Khi nhiệt độ không khí tăng, tốc độ truyền âm trong không khí nhìn chung — phát biểu đúng là: tăng.}
{Khi nhiệt độ không khí tăng, tốc độ truyền âm trong không khí nhìn chung — phát biểu không đổi là đúng.}
\loigiai{
\begin{itemchoice}
\item (Đúng) Trong thí nghiệm, đại lượng thường được đo trực tiếp là — phát biểu đúng là: chiều dài cột khí tại cộng hưởng. (Vì): Ta đọc trực tiếp chiều dài cột khí; bước sóng và tốc độ được tính gián tiếp. b) Sai: phương án nêu không phù hợp. c) Đúng: Tốc độ âm trong khí tăng khi nhiệt độ tăng. d) Sai: phương án nêu không phù hợp
\item (Sai) Trong thí nghiệm, đại lượng thường được đo trực tiếp là — phát biểu tốc độ âm là đúng.
\item (Đúng) Khi nhiệt độ không khí tăng, tốc độ truyền âm trong không khí nhìn chung — phát biểu đúng là: tăng.
\item (Sai) Khi nhiệt độ không khí tăng, tốc độ truyền âm trong không khí nhìn chung — phát biểu không đổi là đúng.
\end{itemchoice}
}
\end{ex}

% ===== Câu 118 =====
% ID: L11C2B15-30-TL
% Mức: VDC
\dangbt{Xử lí số liệu và đánh giá sai số phép đo}
\begin{ex}
% ID: L11C2B15-30-TL
% Mức: VDC
Trong dạng “Xử lí số liệu và đánh giá sai số phép đo”, Một nhóm đo được $342\,\text{m}$ /s trong khi giá trị tham chiếu là $340\,\text{m}$ /s. Nhận xét kết quả.
\choiceTF

\loigiai{
Độ lệch tuyệt đối là $2\,\text{m}$ /s, sai lệch tương đối khoảng $2/340\cdot100\%=0,59\%$; kết quả khá phù hợp trong điều kiện thí nghiệm phổ thông.
}
\end{ex}

% ===== Câu 119 =====
% ID: L11C2B15-30-TLN
% Mức: VDC
\begin{ex}
% ID: L11C2B15-30-TLN
% Mức: VDC
Trong dạng “Xử lí số liệu và đánh giá sai số phép đo”, Đo được $d=0,30$ m với âm $600\,\text{Hz}$. Tính tốc độ âm theo m/s.
\shortans{360}
\loigiai{
$\lambda=2d=0,60$ m; $v=f\lambda=600\cdot0,60=360$ m/s.
}
\end{ex}

% ===== Câu 120 =====
% ID: L11C2B15-30-TN
% Mức: VDC
\dangbt{Tính tốc độ truyền âm từ số liệu thí nghiệm}
\begin{ex}
% ID: L11C2B15-30-TN
% Mức: VDC
Trong dạng “Tính tốc độ truyền âm từ số liệu thí nghiệm”, Khi nhiệt độ không khí tăng, tốc độ truyền âm trong không khí nhìn chung
\choice
{\True tăng}
{giảm}
{không đổi}
{bằng 0}
\loigiai{
Đáp án A. Tốc độ âm trong khí tăng khi nhiệt độ tăng.
}
\end{ex}

% ===== Câu 121 =====
% ID: L11C2B15-31-TL
% Mức: TH
\dangbt{Đề xuất cải tiến và vận dụng thí nghiệm đo tốc độ truyền âm}
\begin{ex}
% ID: L11C2B15-31-TL
% Mức: TH
Trong dạng “Đề xuất cải tiến và vận dụng thí nghiệm đo tốc độ truyền âm”, Nêu các bước chính đo tốc độ truyền âm bằng ống cộng hưởng.
\choiceTF

\loigiai{
Phát âm có tần số f đã biết; thay đổi chiều dài cột khí; ghi hai vị trí cộng hưởng liên tiếp; tính $\lambda=2|L_2-L_1|$ rồi tính $v=f\lambda$.
}
\end{ex}

% ===== Câu 122 =====
% ID: L11C2B15-31-TLN
% Mức: TH
\begin{ex}
% ID: L11C2B15-31-TLN
% Mức: TH
Trong dạng “Đề xuất cải tiến và vận dụng thí nghiệm đo tốc độ truyền âm”, Hai vị trí cộng hưởng liên tiếp cách nhau $34\,\text{cm}$. Tính bước sóng theo cm.
\shortans{68}
\loigiai{
$\lambda=2d=2\cdot34=68$ cm.
}
\end{ex}

% ===== Câu 123 =====
% ID: L11C2B15-31-TN
% Mức: NB
\dangbt{Xác định các vị trí cộng hưởng liên tiếp}
\begin{ex}
% ID: L11C2B15-31-TN
% Mức: NB
Trong dạng “Xác định các vị trí cộng hưởng liên tiếp”, Mục đích chính của thí nghiệm dùng ống cộng hưởng là xác định
\choice
{\True tốc độ truyền âm trong không khí}
{cường độ dòng điện}
{khối lượng riêng của nước}
{nhiệt dung riêng của không khí}
\loigiai{
Đáp án A. Từ tần số đã biết và bước sóng đo qua cộng hưởng, ta tính tốc độ truyền âm.
}
\end{ex}

% ===== Câu 124 =====
% ID: L11C2B15-32-TL
% Mức: TH
\dangbt{Đề xuất cải tiến và vận dụng thí nghiệm đo tốc độ truyền âm}
\begin{ex}
% ID: L11C2B15-32-TL
% Mức: TH
Trong dạng “Đề xuất cải tiến và vận dụng thí nghiệm đo tốc độ truyền âm”, Giải thích vì sao dùng hai vị trí cộng hưởng liên tiếp giúp xác định bước sóng.
\choiceTF

\loigiai{
Trong ống một đầu kín, các chiều dài cộng hưởng kế tiếp chênh nhau $\lambda/2$, nên đo hiệu hai vị trí loại được phần hiệu chỉnh đầu ống và cho $\lambda=2\Delta L$.
}
\end{ex}

% ===== Câu 125 =====
% ID: L11C2B15-32-TLN
% Mức: TH
\begin{ex}
% ID: L11C2B15-32-TLN
% Mức: TH
Trong dạng “Đề xuất cải tiến và vận dụng thí nghiệm đo tốc độ truyền âm”, Âm có tần số $500\,\text{Hz}$, bước sóng $0,68\,\text{m}$. Tính tốc độ theo m/s.
\shortans{340}
\loigiai{
$v=f\lambda=500\cdot0,68=340$ m/s.
}
\end{ex}

% ===== Câu 126 =====
% ID: L11C2B15-32-TN
% Mức: NB
\dangbt{Xác định các vị trí cộng hưởng liên tiếp}
\begin{ex}
% ID: L11C2B15-32-TN
% Mức: NB
Trong dạng “Xác định các vị trí cộng hưởng liên tiếp”, Dụng cụ phát âm có tần số xác định trong thí nghiệm thường là
\choice
{\True âm thoa hoặc loa nối máy phát âm tần}
{nhiệt kế}
{lực kế}
{cân điện tử}
\loigiai{
Đáp án A. Âm thoa hoặc loa tạo âm có tần số biết trước, cần cho phép đo.
}
\end{ex}

% ===== Câu 127 =====
% ID: L11C2B15-33-TL
% Mức: VD
\dangbt{Đề xuất cải tiến và vận dụng thí nghiệm đo tốc độ truyền âm}
\begin{ex}
% ID: L11C2B15-33-TL
% Mức: VD
Trong dạng “Đề xuất cải tiến và vận dụng thí nghiệm đo tốc độ truyền âm”, Với $f=400$ Hz, hai vị trí cộng hưởng là $18\,\text{cm}$ và $60\,\text{cm}$. Tính tốc độ âm.
\choiceTF

\loigiai{
$d=60-18=42\,\mathrm{cm}=0,42$ m; $\lambda=0,84$ m; $v=400\cdot0,84=336$ m/s.
}
\end{ex}

% ===== Câu 128 =====
% ID: L11C2B15-33-TLN
% Mức: VD
\begin{ex}
% ID: L11C2B15-33-TLN
% Mức: VD
Trong dạng “Đề xuất cải tiến và vận dụng thí nghiệm đo tốc độ truyền âm”, Các khoảng cách cộng hưởng đo được là $33,8\,\text{cm}$; $34,0\,\text{cm}$; $34,2\,\text{cm}$. Tính giá trị trung bình theo cm.
\shortans{34}
\loigiai{
$\bar d=(33,8+34,0+34,2)/3=34,0$ cm.
}
\end{ex}

% ===== Câu 129 =====
% ID: L11C2B15-33-TN
% Mức: TH
\dangbt{Xác định các vị trí cộng hưởng liên tiếp}
\begin{ex}
% ID: L11C2B15-33-TN
% Mức: TH
Trong dạng “Xác định các vị trí cộng hưởng liên tiếp”, Khoảng cách giữa hai vị trí cộng hưởng liên tiếp trong ống một đầu kín xấp xỉ
\choice
{\True $\lambda/2$}
{$\lambda$}
{$\lambda/4$}
{$2\lambda$}
\loigiai{
Đáp án A. Hai trạng thái cộng hưởng liên tiếp chênh nhau nửa bước sóng.
}
\end{ex}

% ===== Câu 130 =====
% ID: L11C2B15-34-TL
% Mức: VD
\dangbt{Đề xuất cải tiến và vận dụng thí nghiệm đo tốc độ truyền âm}
\begin{ex}
% ID: L11C2B15-34-TL
% Mức: VD
Trong dạng “Đề xuất cải tiến và vận dụng thí nghiệm đo tốc độ truyền âm”, Nêu hai biện pháp giảm sai số của phép đo.
\choiceTF

\loigiai{
Đặt mắt vuông góc khi đọc thước; xác định nhiều cặp cộng hưởng, đo lặp lại và lấy trung bình; giữ tần số, nhiệt độ ổn định.
}
\end{ex}

% ===== Câu 131 =====
% ID: L11C2B15-34-TLN
% Mức: VD
\begin{ex}
% ID: L11C2B15-34-TLN
% Mức: VD
Trong dạng “Đề xuất cải tiến và vận dụng thí nghiệm đo tốc độ truyền âm”, Với $\bar d=34$ cm và $f=500$ Hz, tính tốc độ âm theo m/s.
\shortans{340}
\loigiai{
$\lambda=2d=0,68$ m; $v=500\cdot0,68=340$ m/s.
}
\end{ex}

% ===== Câu 132 =====
% ID: L11C2B15-34-TN
% Mức: TH
\dangbt{Xác định các vị trí cộng hưởng liên tiếp}
\begin{ex}
% ID: L11C2B15-34-TN
% Mức: TH
Trong dạng “Xác định các vị trí cộng hưởng liên tiếp”, Biết khoảng cách hai vị trí cộng hưởng liên tiếp là d, bước sóng bằng
\choice
{\True $2d$}
{$d/2$}
{$d$}
{$4d$}
\loigiai{
Đáp án A. Vì $d=\lambda/2$ nên $\lambda=2d$.
}
\end{ex}

% ===== Câu 133 =====
% ID: L11C2B15-35-TL
% Mức: VD
\dangbt{Đề xuất cải tiến và vận dụng thí nghiệm đo tốc độ truyền âm}
\begin{ex}
% ID: L11C2B15-35-TL
% Mức: VD
Trong dạng “Đề xuất cải tiến và vận dụng thí nghiệm đo tốc độ truyền âm”, Phân biệt đại lượng đo trực tiếp và gián tiếp trong thí nghiệm.
\choiceTF

\loigiai{
Chiều dài cột khí và tần số hiển thị là số liệu trực tiếp; bước sóng và tốc độ âm được suy ra bằng công thức nên là đại lượng gián tiếp.
}
\end{ex}

% ===== Câu 134 =====
% ID: L11C2B15-35-TLN
% Mức: VD
\begin{ex}
% ID: L11C2B15-35-TLN
% Mức: VD
Trong dạng “Đề xuất cải tiến và vận dụng thí nghiệm đo tốc độ truyền âm”, Giá trị đo là $338\,\text{m}$ /s, giá trị tham chiếu $340\,\text{m}$ /s. Tính độ lệch tuyệt đối theo m/s.
\shortans{2}
\loigiai{
$\Delta v=|338-340|=2$ m/s.
}
\end{ex}

% ===== Câu 135 =====
% ID: L11C2B15-35-TN
% Mức: TH
\dangbt{Xác định các vị trí cộng hưởng liên tiếp}
\begin{ex}
% ID: L11C2B15-35-TN
% Mức: TH
Trong dạng “Xác định các vị trí cộng hưởng liên tiếp”, Công thức tính tốc độ truyền âm là
\choice
{\True $v=f\lambda$}
{$v=f/\lambda$}
{$v=\lambda/f$}
{$v=f+\lambda$}
\loigiai{
Đáp án A. Tốc độ truyền sóng thỏa $v=f\lambda$.
}
\end{ex}

% ===== Câu 136 =====
% ID: L11C2B15-36-TL
% Mức: VDC
\dangbt{Đề xuất cải tiến và vận dụng thí nghiệm đo tốc độ truyền âm}
\begin{ex}
% ID: L11C2B15-36-TL
% Mức: VDC
Trong dạng “Đề xuất cải tiến và vận dụng thí nghiệm đo tốc độ truyền âm”, Một nhóm đo được $342\,\text{m}$ /s trong khi giá trị tham chiếu là $340\,\text{m}$ /s. Nhận xét kết quả.
\choiceTF

\loigiai{
Độ lệch tuyệt đối là $2\,\text{m}$ /s, sai lệch tương đối khoảng $2/340\cdot100\%=0,59\%$; kết quả khá phù hợp trong điều kiện thí nghiệm phổ thông.
}
\end{ex}

% ===== Câu 137 =====
% ID: L11C2B15-36-TLN
% Mức: VDC
\begin{ex}
% ID: L11C2B15-36-TLN
% Mức: VDC
Trong dạng “Đề xuất cải tiến và vận dụng thí nghiệm đo tốc độ truyền âm”, Đo được $d=0,30$ m với âm $600\,\text{Hz}$. Tính tốc độ âm theo m/s.
\shortans{360}
\loigiai{
$\lambda=2d=0,60$ m; $v=f\lambda=600\cdot0,60=360$ m/s.
}
\end{ex}

% ===== Câu 138 =====
% ID: L11C2B15-36-TN
% Mức: VD
\dangbt{Xác định các vị trí cộng hưởng liên tiếp}
\begin{ex}
% ID: L11C2B15-36-TN
% Mức: VD
Trong dạng “Xác định các vị trí cộng hưởng liên tiếp”, Để giảm sai số khi xác định khoảng cách cộng hưởng nên
\choice
{\True đo nhiều cặp vị trí và lấy trung bình}
{chỉ đo một lần}
{thay đổi tần số liên tục trong mỗi lần đo}
{đọc thước từ góc xiên}
\loigiai{
Đáp án A. Đo lặp lại và lấy trung bình làm giảm ảnh hưởng sai số ngẫu nhiên.
}
\end{ex}

% ===== Câu 139 =====
% ID: L11C2B15-37-TN
% Mức: VD
\begin{ex}
% ID: L11C2B15-37-TN
% Mức: VD
Trong dạng “Xác định các vị trí cộng hưởng liên tiếp”, Khi đọc thước, mắt phải đặt
\choice
{\True vuông góc với vạch chia tại vị trí cần đọc}
{càng xa thước càng tốt}
{xiên 45 độ}
{ở bất kì vị trí nào}
\loigiai{
Đáp án A. Đặt mắt vuông góc giúp tránh sai số thị sai.
}
\end{ex}

% ===== Câu 140 =====
% ID: L11C2B15-38-TN
% Mức: VD
\begin{ex}
% ID: L11C2B15-38-TN
% Mức: VD
Trong dạng “Xác định các vị trí cộng hưởng liên tiếp”, Nếu $f=500$ Hz và $\lambda=0,68$ m thì tốc độ âm là
\choice
{\True $340\,\text{m}$ /s}
{$735\,\text{m}$ /s}
{$0,00136\,\text{m}$ /s}
{$500,68\,\text{m}$ /s}
\loigiai{
Đáp án A. $v=f\lambda=500\cdot0,68=340$ m/s.
}
\end{ex}

% ===== Câu 141 =====
% ID: L11C2B15-39-TN
% Mức: VDC
\begin{ex}
% ID: L11C2B15-39-TN
% Mức: VDC
Trong dạng “Xác định các vị trí cộng hưởng liên tiếp”, Trong thí nghiệm, đại lượng thường được đo trực tiếp là
\choice
{\True chiều dài cột khí tại cộng hưởng}
{tốc độ âm}
{bước sóng chính xác tuyệt đối}
{pha ban đầu của nguồn}
\loigiai{
Đáp án A. Ta đọc trực tiếp chiều dài cột khí; bước sóng và tốc độ được tính gián tiếp.
}
\end{ex}

% ===== Câu 142 =====
% ID: L11C2B15-40-TN
% Mức: VDC
\begin{ex}
% ID: L11C2B15-40-TN
% Mức: VDC
Trong dạng “Xác định các vị trí cộng hưởng liên tiếp”, Khi nhiệt độ không khí tăng, tốc độ truyền âm trong không khí nhìn chung
\choice
{\True tăng}
{giảm}
{không đổi}
{bằng 0}
\loigiai{
Đáp án A. Tốc độ âm trong khí tăng khi nhiệt độ tăng.
}
\end{ex}

% ===== Câu 143 =====
% ID: L11C2B15-41-TN
% Mức: NB
\dangbt{Xử lí số liệu và đánh giá sai số phép đo}
\begin{ex}
% ID: L11C2B15-41-TN
% Mức: NB
Trong dạng “Xử lí số liệu và đánh giá sai số phép đo”, Mục đích chính của thí nghiệm dùng ống cộng hưởng là xác định
\choice
{\True tốc độ truyền âm trong không khí}
{cường độ dòng điện}
{khối lượng riêng của nước}
{nhiệt dung riêng của không khí}
\loigiai{
Đáp án A. Từ tần số đã biết và bước sóng đo qua cộng hưởng, ta tính tốc độ truyền âm.
}
\end{ex}

% ===== Câu 144 =====
% ID: L11C2B15-42-TN
% Mức: NB
\begin{ex}
% ID: L11C2B15-42-TN
% Mức: NB
Trong dạng “Xử lí số liệu và đánh giá sai số phép đo”, Dụng cụ phát âm có tần số xác định trong thí nghiệm thường là
\choice
{\True âm thoa hoặc loa nối máy phát âm tần}
{nhiệt kế}
{lực kế}
{cân điện tử}
\loigiai{
Đáp án A. Âm thoa hoặc loa tạo âm có tần số biết trước, cần cho phép đo.
}
\end{ex}

% ===== Câu 145 =====
% ID: L11C2B15-43-TN
% Mức: TH
\begin{ex}
% ID: L11C2B15-43-TN
% Mức: TH
Trong dạng “Xử lí số liệu và đánh giá sai số phép đo”, Khoảng cách giữa hai vị trí cộng hưởng liên tiếp trong ống một đầu kín xấp xỉ
\choice
{\True $\lambda/2$}
{$\lambda$}
{$\lambda/4$}
{$2\lambda$}
\loigiai{
Đáp án A. Hai trạng thái cộng hưởng liên tiếp chênh nhau nửa bước sóng.
}
\end{ex}

% ===== Câu 146 =====
% ID: L11C2B15-44-TN
% Mức: TH
\begin{ex}
% ID: L11C2B15-44-TN
% Mức: TH
Trong dạng “Xử lí số liệu và đánh giá sai số phép đo”, Biết khoảng cách hai vị trí cộng hưởng liên tiếp là d, bước sóng bằng
\choice
{\True $2d$}
{$d/2$}
{$d$}
{$4d$}
\loigiai{
Đáp án A. Vì $d=\lambda/2$ nên $\lambda=2d$.
}
\end{ex}

% ===== Câu 147 =====
% ID: L11C2B15-45-TN
% Mức: TH
\begin{ex}
% ID: L11C2B15-45-TN
% Mức: TH
Trong dạng “Xử lí số liệu và đánh giá sai số phép đo”, Công thức tính tốc độ truyền âm là
\choice
{\True $v=f\lambda$}
{$v=f/\lambda$}
{$v=\lambda/f$}
{$v=f+\lambda$}
\loigiai{
Đáp án A. Tốc độ truyền sóng thỏa $v=f\lambda$.
}
\end{ex}

% ===== Câu 148 =====
% ID: L11C2B15-46-TN
% Mức: VD
\begin{ex}
% ID: L11C2B15-46-TN
% Mức: VD
Trong dạng “Xử lí số liệu và đánh giá sai số phép đo”, Để giảm sai số khi xác định khoảng cách cộng hưởng nên
\choice
{\True đo nhiều cặp vị trí và lấy trung bình}
{chỉ đo một lần}
{thay đổi tần số liên tục trong mỗi lần đo}
{đọc thước từ góc xiên}
\loigiai{
Đáp án A. Đo lặp lại và lấy trung bình làm giảm ảnh hưởng sai số ngẫu nhiên.
}
\end{ex}

% ===== Câu 149 =====
% ID: L11C2B15-47-TN
% Mức: VD
\begin{ex}
% ID: L11C2B15-47-TN
% Mức: VD
Trong dạng “Xử lí số liệu và đánh giá sai số phép đo”, Khi đọc thước, mắt phải đặt
\choice
{\True vuông góc với vạch chia tại vị trí cần đọc}
{càng xa thước càng tốt}
{xiên 45 độ}
{ở bất kì vị trí nào}
\loigiai{
Đáp án A. Đặt mắt vuông góc giúp tránh sai số thị sai.
}
\end{ex}

% ===== Câu 150 =====
% ID: L11C2B15-48-TN
% Mức: VD
\begin{ex}
% ID: L11C2B15-48-TN
% Mức: VD
Trong dạng “Xử lí số liệu và đánh giá sai số phép đo”, Nếu $f=500$ Hz và $\lambda=0,68$ m thì tốc độ âm là
\choice
{\True $340\,\text{m}$ /s}
{$735\,\text{m}$ /s}
{$0,00136\,\text{m}$ /s}
{$500,68\,\text{m}$ /s}
\loigiai{
Đáp án A. $v=f\lambda=500\cdot0,68=340$ m/s.
}
\end{ex}

% ===== Câu 151 =====
% ID: L11C2B15-49-TN
% Mức: VDC
\begin{ex}
% ID: L11C2B15-49-TN
% Mức: VDC
Trong dạng “Xử lí số liệu và đánh giá sai số phép đo”, Trong thí nghiệm, đại lượng thường được đo trực tiếp là
\choice
{\True chiều dài cột khí tại cộng hưởng}
{tốc độ âm}
{bước sóng chính xác tuyệt đối}
{pha ban đầu của nguồn}
\loigiai{
Đáp án A. Ta đọc trực tiếp chiều dài cột khí; bước sóng và tốc độ được tính gián tiếp.
}
\end{ex}

% ===== Câu 152 =====
% ID: L11C2B15-50-TN
% Mức: VDC
\begin{ex}
% ID: L11C2B15-50-TN
% Mức: VDC
Trong dạng “Xử lí số liệu và đánh giá sai số phép đo”, Khi nhiệt độ không khí tăng, tốc độ truyền âm trong không khí nhìn chung
\choice
{\True tăng}
{giảm}
{không đổi}
{bằng 0}
\loigiai{
Đáp án A. Tốc độ âm trong khí tăng khi nhiệt độ tăng.
}
\end{ex}

% ===== Câu 153 =====
% ID: L11C2B15-51-TN
% Mức: NB
\dangbt{Đề xuất cải tiến và vận dụng thí nghiệm đo tốc độ truyền âm}
\begin{ex}
% ID: L11C2B15-51-TN
% Mức: NB
Trong dạng “Đề xuất cải tiến và vận dụng thí nghiệm đo tốc độ truyền âm”, Mục đích chính của thí nghiệm dùng ống cộng hưởng là xác định
\choice
{\True tốc độ truyền âm trong không khí}
{cường độ dòng điện}
{khối lượng riêng của nước}
{nhiệt dung riêng của không khí}
\loigiai{
Đáp án A. Từ tần số đã biết và bước sóng đo qua cộng hưởng, ta tính tốc độ truyền âm.
}
\end{ex}

% ===== Câu 154 =====
% ID: L11C2B15-52-TN
% Mức: NB
\begin{ex}
% ID: L11C2B15-52-TN
% Mức: NB
Trong dạng “Đề xuất cải tiến và vận dụng thí nghiệm đo tốc độ truyền âm”, Dụng cụ phát âm có tần số xác định trong thí nghiệm thường là
\choice
{\True âm thoa hoặc loa nối máy phát âm tần}
{nhiệt kế}
{lực kế}
{cân điện tử}
\loigiai{
Đáp án A. Âm thoa hoặc loa tạo âm có tần số biết trước, cần cho phép đo.
}
\end{ex}

% ===== Câu 155 =====
% ID: L11C2B15-53-TN
% Mức: TH
\begin{ex}
% ID: L11C2B15-53-TN
% Mức: TH
Trong dạng “Đề xuất cải tiến và vận dụng thí nghiệm đo tốc độ truyền âm”, Khoảng cách giữa hai vị trí cộng hưởng liên tiếp trong ống một đầu kín xấp xỉ
\choice
{\True $\lambda/2$}
{$\lambda$}
{$\lambda/4$}
{$2\lambda$}
\loigiai{
Đáp án A. Hai trạng thái cộng hưởng liên tiếp chênh nhau nửa bước sóng.
}
\end{ex}

% ===== Câu 156 =====
% ID: L11C2B15-54-TN
% Mức: TH
\begin{ex}
% ID: L11C2B15-54-TN
% Mức: TH
Trong dạng “Đề xuất cải tiến và vận dụng thí nghiệm đo tốc độ truyền âm”, Biết khoảng cách hai vị trí cộng hưởng liên tiếp là d, bước sóng bằng
\choice
{\True $2d$}
{$d/2$}
{$d$}
{$4d$}
\loigiai{
Đáp án A. Vì $d=\lambda/2$ nên $\lambda=2d$.
}
\end{ex}

% ===== Câu 157 =====
% ID: L11C2B15-55-TN
% Mức: TH
\begin{ex}
% ID: L11C2B15-55-TN
% Mức: TH
Trong dạng “Đề xuất cải tiến và vận dụng thí nghiệm đo tốc độ truyền âm”, Công thức tính tốc độ truyền âm là
\choice
{\True $v=f\lambda$}
{$v=f/\lambda$}
{$v=\lambda/f$}
{$v=f+\lambda$}
\loigiai{
Đáp án A. Tốc độ truyền sóng thỏa $v=f\lambda$.
}
\end{ex}

% ===== Câu 158 =====
% ID: L11C2B15-56-TN
% Mức: VD
\begin{ex}
% ID: L11C2B15-56-TN
% Mức: VD
Trong dạng “Đề xuất cải tiến và vận dụng thí nghiệm đo tốc độ truyền âm”, Để giảm sai số khi xác định khoảng cách cộng hưởng nên
\choice
{\True đo nhiều cặp vị trí và lấy trung bình}
{chỉ đo một lần}
{thay đổi tần số liên tục trong mỗi lần đo}
{đọc thước từ góc xiên}
\loigiai{
Đáp án A. Đo lặp lại và lấy trung bình làm giảm ảnh hưởng sai số ngẫu nhiên.
}
\end{ex}

% ===== Câu 159 =====
% ID: L11C2B15-57-TN
% Mức: VD
\begin{ex}
% ID: L11C2B15-57-TN
% Mức: VD
Trong dạng “Đề xuất cải tiến và vận dụng thí nghiệm đo tốc độ truyền âm”, Khi đọc thước, mắt phải đặt
\choice
{\True vuông góc với vạch chia tại vị trí cần đọc}
{càng xa thước càng tốt}
{xiên 45 độ}
{ở bất kì vị trí nào}
\loigiai{
Đáp án A. Đặt mắt vuông góc giúp tránh sai số thị sai.
}
\end{ex}

% ===== Câu 160 =====
% ID: L11C2B15-58-TN
% Mức: VD
\begin{ex}
% ID: L11C2B15-58-TN
% Mức: VD
Trong dạng “Đề xuất cải tiến và vận dụng thí nghiệm đo tốc độ truyền âm”, Nếu $f=500$ Hz và $\lambda=0,68$ m thì tốc độ âm là
\choice
{\True $340\,\text{m}$ /s}
{$735\,\text{m}$ /s}
{$0,00136\,\text{m}$ /s}
{$500,68\,\text{m}$ /s}
\loigiai{
Đáp án A. $v=f\lambda=500\cdot0,68=340$ m/s.
}
\end{ex}

% ===== Câu 161 =====
% ID: L11C2B15-59-TN
% Mức: VDC
\begin{ex}
% ID: L11C2B15-59-TN
% Mức: VDC
Trong dạng “Đề xuất cải tiến và vận dụng thí nghiệm đo tốc độ truyền âm”, Trong thí nghiệm, đại lượng thường được đo trực tiếp là
\choice
{\True chiều dài cột khí tại cộng hưởng}
{tốc độ âm}
{bước sóng chính xác tuyệt đối}
{pha ban đầu của nguồn}
\loigiai{
Đáp án A. Ta đọc trực tiếp chiều dài cột khí; bước sóng và tốc độ được tính gián tiếp.
}
\end{ex}

% ===== Câu 162 =====
% ID: L11C2B15-60-TN
% Mức: VDC
\begin{ex}
% ID: L11C2B15-60-TN
% Mức: VDC
Trong dạng “Đề xuất cải tiến và vận dụng thí nghiệm đo tốc độ truyền âm”, Khi nhiệt độ không khí tăng, tốc độ truyền âm trong không khí nhìn chung
\choice
{\True tăng}
{giảm}
{không đổi}
{bằng 0}
\loigiai{
Đáp án A. Tốc độ âm trong khí tăng khi nhiệt độ tăng.
}
\end{ex}
